\documentclass[12pt]{article}
\usepackage[utf8]{inputenc}
\usepackage{amsmath}
\usepackage{amssymb}
\usepackage{amsthm}
\usepackage{graphicx}
\usepackage{hyperref}
\usepackage{geometry}
\usepackage{fancyhdr}
\usepackage{listings}
\usepackage{color}

\geometry{a4paper, margin=1in}

\title{超大型 LaTeX 文档测试}
\author{Code Reader 性能测试}
\date{\today}

\newtheorem{theorem}{定理}[section]
\newtheorem{lemma}[theorem]{引理}
\newtheorem{proposition}[theorem]{命题}
\newtheorem{corollary}[theorem]{推论}
\newtheorem{definition}{定义}[section]
\newtheorem{example}{例}[section]
\newtheorem{remark}{注}[section]

\begin{document}

\maketitle

\begin{abstract}
本文档是一个超大型 LaTeX 测试文档，用于测试 Code Reader 的大文件处理能力、Tree-sitter 增量解析性能以及 Web Worker 渲染效率。文档包含大量数学公式、定理证明、代码块和复杂排版元素，总大小超过 200KB。
\end{abstract}

\tableofcontents

\newpage

% ================================
% 第一章：数学分析
% ================================

\section{数学分析基础}

\subsection{实数系与极限理论}

\begin{definition}[实数的完备性]
实数集 $\mathbb{R}$ 具有完备性，即每一个柯西序列都收敛于 $\mathbb{R}$ 中的某个元素。形式化地，对于序列 $\{a_n\}$，如果对于任意 $\epsilon > 0$，存在 $N \in \mathbb{N}$，使得当 $m, n > N$ 时，$|a_m - a_n| < \epsilon$，则存在 $a \in \mathbb{R}$ 使得 $\lim_{n \to \infty} a_n = a$。
\end{definition}

\begin{theorem}[确界原理]
设 $S$ 是 $\mathbb{R}$ 的非空子集。如果 $S$ 有上界，则 $S$ 必有上确界；如果 $S$ 有下界，则 $S$ 必有下确界。即：
\begin{align}
\exists M \in \mathbb{R}, \forall x \in S: x \leq M &\Rightarrow \exists \sup S \in \mathbb{R} \\
\exists m \in \mathbb{R}, \forall x \in S: x \geq m &\Rightarrow \exists \inf S \in \mathbb{R}
\end{align}
\end{theorem}

\begin{proof}
设 $S$ 是有上界的非空子集。令 $A = \{x \in \mathbb{R} : x \text{ 是 } S \text{ 的上界}\}$，$B = \{x \in \mathbb{R} : x \text{ 不是 } S \text{ 的上界}\}$。

由于 $S$ 非空且有上界，$A$ 和 $B$ 都非空。显然 $A \cup B = \mathbb{R}$，$A \cap B = \emptyset$。

对于任意 $a \in A$ 和 $b \in B$，必有 $b < a$（否则存在 $s \in S$ 使得 $s > a$，矛盾）。

由戴德金分割原理，存在唯一的实数 $c$ 使得对于所有 $b \in B$ 和 $a \in A$，有 $b \leq c \leq a$。

我们证明 $c = \sup S$：
\begin{enumerate}
\item $c$ 是 $S$ 的上界：若存在 $s \in S$ 使得 $s > c$，取 $\epsilon = s - c$，则 $c + \epsilon/2$ 不是上界，故 $c + \epsilon/2 \in B$，但 $c + \epsilon/2 > c$，矛盾。
\item $c$ 是最小上界：若 $c' < c$ 是上界，则 $c' \in A$，但 $c' < c$ 与 $c$ 的定义矛盾。
\end{enumerate}
\end{proof}

\begin{theorem}[单调收敛定理]
单调有界数列必收敛。具体地：
\begin{enumerate}
\item 若 $\{a_n\}$ 单调递增且有上界，则 $\lim_{n \to \infty} a_n = \sup\{a_n : n \in \mathbb{N}\}$
\item 若 $\{a_n\}$ 单调递减且有下界，则 $\lim_{n \to \infty} a_n = \inf\{a_n : n \in \mathbb{N}\}$
\end{enumerate}
\end{theorem}

\begin{theorem}[区间套定理]
设 $\{[a_n, b_n]\}$ 是一列闭区间，满足：
\begin{enumerate}
\item $[a_{n+1}, b_{n+1}] \subseteq [a_n, b_n]$ 对所有 $n \in \mathbb{N}$ 成立
\item $\lim_{n \to \infty} (b_n - a_n) = 0$
\end{enumerate}
则存在唯一的 $\xi \in \mathbb{R}$ 使得 $\xi \in \bigcap_{n=1}^{\infty} [a_n, b_n]$，即 $\{\xi\} = \bigcap_{n=1}^{\infty} [a_n, b_n]$。
\end{theorem}

\begin{theorem}[Bolzano-Weierstrass 定理]
有界数列必有收敛子列。形式化地，若 $\{a_n\}$ 有界，则存在严格递增的正整数序列 $\{n_k\}$ 和实数 $a$，使得 $\lim_{k \to \infty} a_{n_k} = a$。
\end{theorem}

\begin{proof}
设 $\{a_n\}$ 有界，即存在 $M > 0$ 使得 $|a_n| \leq M$ 对所有 $n$ 成立。

使用区间套方法构造收敛子列：

令 $[c_0, d_0] = [-M, M]$。将 $[c_0, d_0]$ 二等分为 $[c_0, \frac{c_0+d_0}{2}]$ 和 $[\frac{c_0+d_0}{2}, d_0]$。

至少有一个子区间包含 $\{a_n\}$ 的无穷多项，选择这样的子区间作为 $[c_1, d_1]$。

重复此过程，得到区间套 $\{[c_n, d_n]\}$，满足：
\begin{enumerate}
\item $[c_{n+1}, d_{n+1}] \subseteq [c_n, d_n]$
\item $d_n - c_n = \frac{2M}{2^n}$
\item 每个 $[c_n, d_n]$ 包含 $\{a_k\}$ 的无穷多项
\end{enumerate}

由区间套定理，存在唯一的 $\xi = \bigcap_{n=1}^{\infty} [c_n, d_n]$。

构造子列：选 $n_1$ 使 $a_{n_1} \in [c_1, d_1]$，选 $n_2 > n_1$ 使 $a_{n_2} \in [c_2, d_2]$，依此类推。

对于 $\epsilon > 0$，取 $N$ 使 $\frac{2M}{2^N} < \epsilon$，则当 $k > N$ 时，$|a_{n_k} - \xi| < \frac{2M}{2^k} < \epsilon$。

因此 $\lim_{k \to \infty} a_{n_k} = \xi$。
\end{proof}

\subsection{连续函数的性质}

\begin{definition}[函数连续性]
设函数 $f: D \to \mathbb{R}$，$x_0 \in D$。若对于任意 $\epsilon > 0$，存在 $\delta > 0$，使得当 $x \in D$ 且 $|x - x_0| < \delta$ 时，有 $|f(x) - f(x_0)| < \epsilon$，则称 $f$ 在 $x_0$ 处连续。

等价地，$f$ 在 $x_0$ 处连续当且仅当 $\lim_{x \to x_0} f(x) = f(x_0)$。
\end{definition}

\begin{definition}[一致连续性]
设函数 $f: D \to \mathbb{R}$。若对于任意 $\epsilon > 0$，存在 $\delta > 0$，使得对于所有满足 $|x - y| < \delta$ 的 $x, y \in D$，有 $|f(x) - f(y)| < \epsilon$，则称 $f$ 在 $D$ 上一致连续。
\end{definition}

\begin{theorem}[介值定理]
设 $f$ 在闭区间 $[a, b]$ 上连续，$f(a) \neq f(b)$。若 $c$ 是介于 $f(a)$ 和 $f(b)$ 之间的任意实数，则存在 $\xi \in (a, b)$ 使得 $f(\xi) = c$。
\end{theorem}

\begin{proof}
不失一般性，设 $f(a) < c < f(b)$。令 $S = \{x \in [a, b] : f(x) < c\}$。

$S$ 非空（因为 $a \in S$）且有上界 $b$，由确界原理，存在 $\xi = \sup S$。

我们证明 $f(\xi) = c$：

假设 $f(\xi) < c$。由连续性，存在 $\delta > 0$ 使得当 $|x - \xi| < \delta$ 时，$f(x) < c$。取 $x = \xi + \delta/2 \in S$，与 $\xi = \sup S$ 矛盾。

假设 $f(\xi) > c$。由连续性，存在 $\delta > 0$ 使得当 $|x - \xi| < \delta$ 时，$f(x) > c$。则 $\xi - \delta/2$ 是 $S$ 的上界，与 $\xi = \sup S$ 矛盾。

因此 $f(\xi) = c$。
\end{proof}

\begin{theorem}[最值定理]
设 $f$ 在闭区间 $[a, b]$ 上连续，则 $f$ 在 $[a, b]$ 上取得最大值和最小值。即存在 $x_M, x_m \in [a, b]$ 使得对于所有 $x \in [a, b]$：
$$f(x_m) \leq f(x) \leq f(x_M)$$
\end{theorem}

\begin{theorem}[Heine-Cantor 定理]
若 $f$ 在闭区间 $[a, b]$ 上连续，则 $f$ 在 $[a, b]$ 上一致连续。
\end{theorem}

\begin{proof}
反证法。假设 $f$ 在 $[a, b]$ 上不一致连续，则存在 $\epsilon_0 > 0$，使得对于任意 $\delta > 0$，存在 $x, y \in [a, b]$ 满足 $|x - y| < \delta$ 但 $|f(x) - f(y)| \geq \epsilon_0$。

取 $\delta = 1/n$，得到序列 $\{x_n\}$ 和 $\{y_n\}$，满足：
\begin{enumerate}
\item $|x_n - y_n| < 1/n$
\item $|f(x_n) - f(y_n)| \geq \epsilon_0$
\end{enumerate}

由 Bolzano-Weierstrass 定理，$\{x_n\}$ 有收敛子列 $\{x_{n_k}\}$，设 $\lim_{k \to \infty} x_{n_k} = \xi$。

由 $|x_{n_k} - y_{n_k}| < 1/n_k \to 0$，有 $\lim_{k \to \infty} y_{n_k} = \xi$。

由 $f$ 在 $\xi$ 处连续：
$$\lim_{k \to \infty} f(x_{n_k}) = f(\xi) = \lim_{k \to \infty} f(y_{n_k})$$

因此 $\lim_{k \to \infty} |f(x_{n_k}) - f(y_{n_k})| = 0$，与 $|f(x_{n_k}) - f(y_{n_k})| \geq \epsilon_0$ 矛盾。
\end{proof}

\subsection{微分学}

\begin{definition}[导数]
设函数 $f$ 在 $x_0$ 的某邻域内有定义。若极限
$$f'(x_0) = \lim_{h \to 0} \frac{f(x_0 + h) - f(x_0)}{h}$$
存在，则称 $f$ 在 $x_0$ 处可导，此极限值称为 $f$ 在 $x_0$ 处的导数。
\end{definition}

\begin{theorem}[可导必连续]
若 $f$ 在 $x_0$ 处可导，则 $f$ 在 $x_0$ 处连续。
\end{theorem}

\begin{proof}
设 $f'(x_0)$ 存在。则：
$$\lim_{x \to x_0} [f(x) - f(x_0)] = \lim_{x \to x_0} \frac{f(x) - f(x_0)}{x - x_0} \cdot (x - x_0) = f'(x_0) \cdot 0 = 0$$
因此 $\lim_{x \to x_0} f(x) = f(x_0)$，即 $f$ 在 $x_0$ 处连续。
\end{proof}

\begin{theorem}[求导法则]
设 $f$ 和 $g$ 在 $x$ 处可导，则：
\begin{enumerate}
\item 和差法则：$(f \pm g)'(x) = f'(x) \pm g'(x)$
\item 积法则：$(fg)'(x) = f'(x)g(x) + f(x)g'(x)$
\item 商法则（$g(x) \neq 0$）：$\left(\frac{f}{g}\right)'(x) = \frac{f'(x)g(x) - f(x)g'(x)}{[g(x)]^2}$
\end{enumerate}
\end{theorem}

\begin{theorem}[链式法则]
设 $g$ 在 $x$ 处可导，$f$ 在 $g(x)$ 处可导，则复合函数 $f \circ g$ 在 $x$ 处可导，且：
$$(f \circ g)'(x) = f'(g(x)) \cdot g'(x)$$
\end{theorem}

\begin{theorem}[反函数求导]
设 $f$ 在区间 $I$ 上严格单调且连续，$f'(x_0) \neq 0$，则 $f^{-1}$ 在 $y_0 = f(x_0)$ 处可导，且：
$$(f^{-1})'(y_0) = \frac{1}{f'(x_0)} = \frac{1}{f'(f^{-1}(y_0))}$$
\end{theorem}

\begin{theorem}[费马引理]
设 $f$ 在 $x_0$ 处取得极值，且 $f$ 在 $x_0$ 处可导，则 $f'(x_0) = 0$。
\end{theorem}

\begin{theorem}[罗尔定理]
设 $f$ 在闭区间 $[a, b]$ 上连续，在开区间 $(a, b)$ 内可导，且 $f(a) = f(b)$。则存在 $\xi \in (a, b)$ 使得 $f'(\xi) = 0$。
\end{theorem}

\begin{proof}
若 $f$ 是常函数，则对任意 $\xi \in (a, b)$ 有 $f'(\xi) = 0$。

若 $f$ 不是常函数，由最值定理，$f$ 在 $[a, b]$ 上取得最大值 $M$ 和最小值 $m$。

由 $f(a) = f(b)$ 且 $f$ 不是常函数，最大值或最小值在 $(a, b)$ 内取得。

设 $f(\xi) = M$（或 $m$），$\xi \in (a, b)$。由费马引理，$f'(\xi) = 0$。
\end{proof}

\begin{theorem}[拉格朗日中值定理]
设 $f$ 在闭区间 $[a, b]$ 上连续，在开区间 $(a, b)$ 内可导。则存在 $\xi \in (a, b)$ 使得：
$$f'(\xi) = \frac{f(b) - f(a)}{b - a}$$
\end{theorem}

\begin{proof}
构造辅助函数：
$$g(x) = f(x) - \frac{f(b) - f(a)}{b - a}(x - a)$$

则 $g$ 满足罗尔定理条件：
\begin{enumerate}
\item $g$ 在 $[a, b]$ 上连续
\item $g$ 在 $(a, b)$ 内可导
\item $g(a) = f(a)$，$g(b) = f(b) - [f(b) - f(a)] = f(a) = g(a)$
\end{enumerate}

由罗尔定理，存在 $\xi \in (a, b)$ 使得 $g'(\xi) = 0$，即：
$$f'(\xi) - \frac{f(b) - f(a)}{b - a} = 0$$
因此 $f'(\xi) = \frac{f(b) - f(a)}{b - a}$。
\end{proof}

\begin{theorem}[柯西中值定理]
设 $f, g$ 在闭区间 $[a, b]$ 上连续，在开区间 $(a, b)$ 内可导，且 $g'(x) \neq 0$ 对所有 $x \in (a, b)$ 成立。则存在 $\xi \in (a, b)$ 使得：
$$\frac{f'(\xi)}{g'(\xi)} = \frac{f(b) - f(a)}{g(b) - g(a)}$$
\end{theorem}

\begin{theorem}[洛必达法则]
设 $\lim_{x \to a} f(x) = \lim_{x \to a} g(x) = 0$（或 $\pm\infty$），$f$ 和 $g$ 在 $a$ 的去心邻域内可导，$g'(x) \neq 0$。若 $\lim_{x \to a} \frac{f'(x)}{g'(x)}$ 存在（或为 $\pm\infty$），则：
$$\lim_{x \to a} \frac{f(x)}{g(x)} = \lim_{x \to a} \frac{f'(x)}{g'(x)}$$
\end{theorem}

\begin{theorem}[泰勒公式]
设 $f$ 在 $x_0$ 处 $n$ 阶可导，则：
$$f(x) = \sum_{k=0}^{n} \frac{f^{(k)}(x_0)}{k!}(x - x_0)^k + o((x - x_0)^n)$$
这称为 $f$ 在 $x_0$ 处的 $n$ 阶泰勒展开（带 Peano 余项）。

若 $f$ 在包含 $x_0$ 的某区间上 $n+1$ 阶可导，则余项可写为 Lagrange 形式：
$$R_n(x) = \frac{f^{(n+1)}(\xi)}{(n+1)!}(x - x_0)^{n+1}$$
其中 $\xi$ 在 $x$ 和 $x_0$ 之间。
\end{theorem}

常用泰勒展开（$x_0 = 0$，即麦克劳林展开）：
\begin{align}
e^x &= \sum_{n=0}^{\infty} \frac{x^n}{n!} = 1 + x + \frac{x^2}{2!} + \frac{x^3}{3!} + \cdots \\
\sin x &= \sum_{n=0}^{\infty} \frac{(-1)^n x^{2n+1}}{(2n+1)!} = x - \frac{x^3}{3!} + \frac{x^5}{5!} - \cdots \\
\cos x &= \sum_{n=0}^{\infty} \frac{(-1)^n x^{2n}}{(2n)!} = 1 - \frac{x^2}{2!} + \frac{x^4}{4!} - \cdots \\
\ln(1+x) &= \sum_{n=1}^{\infty} \frac{(-1)^{n+1} x^n}{n} = x - \frac{x^2}{2} + \frac{x^3}{3} - \cdots \quad (-1 < x \leq 1) \\
(1+x)^\alpha &= \sum_{n=0}^{\infty} \binom{\alpha}{n} x^n = 1 + \alpha x + \frac{\alpha(\alpha-1)}{2!}x^2 + \cdots \quad (|x| < 1)
\end{align}

\subsection{积分学}

\begin{definition}[黎曼积分]
设 $f$ 在 $[a, b]$ 上有界。对 $[a, b]$ 的分割 $P: a = x_0 < x_1 < \cdots < x_n = b$，定义：
\begin{align}
m_i &= \inf\{f(x) : x \in [x_{i-1}, x_i]\} \\
M_i &= \sup\{f(x) : x \in [x_{i-1}, x_i]\} \\
L(P, f) &= \sum_{i=1}^{n} m_i \Delta x_i \quad \text{（下和）} \\
U(P, f) &= \sum_{i=1}^{n} M_i \Delta x_i \quad \text{（上和）}
\end{align}

定义下积分和上积分：
\begin{align}
\underline{\int_a^b} f(x) dx &= \sup_P L(P, f) \\
\overline{\int_a^b} f(x) dx &= \inf_P U(P, f)
\end{align}

若下积分等于上积分，则称 $f$ 在 $[a, b]$ 上黎曼可积，公共值记为 $\int_a^b f(x) dx$。
\end{definition}

\begin{theorem}[可积的充要条件]
$f$ 在 $[a, b]$ 上黎曼可积的充要条件是：对于任意 $\epsilon > 0$，存在分割 $P$ 使得 $U(P, f) - L(P, f) < \epsilon$。
\end{theorem}

\begin{theorem}[连续函数可积]
若 $f$ 在 $[a, b]$ 上连续，则 $f$ 在 $[a, b]$ 上黎曼可积。
\end{theorem}

\begin{theorem}[单调函数可积]
若 $f$ 在 $[a, b]$ 上单调，则 $f$ 在 $[a, b]$ 上黎曼可积。
\end{theorem}

\begin{theorem}[积分中值定理]
设 $f$ 在 $[a, b]$ 上连续，则存在 $\xi \in [a, b]$ 使得：
$$\int_a^b f(x) dx = f(\xi)(b - a)$$
\end{theorem}

\begin{theorem}[微积分基本定理（Newton-Leibniz 公式）]
设 $f$ 在 $[a, b]$ 上连续，$F$ 是 $f$ 的一个原函数，则：
$$\int_a^b f(x) dx = F(b) - F(a)$$
\end{theorem}

\begin{theorem}[分部积分公式]
若 $u, v$ 在 $[a, b]$ 上有连续导数，则：
$$\int_a^b u(x) v'(x) dx = [u(x)v(x)]_a^b - \int_a^b u'(x) v(x) dx$$
\end{theorem}

\begin{theorem}[换元积分公式]
设 $f$ 在 $[a, b]$ 上连续，$\varphi: [\alpha, \beta] \to [a, b]$ 连续可微，$\varphi(\alpha) = a$，$\varphi(\beta) = b$，则：
$$\int_a^b f(x) dx = \int_\alpha^\beta f(\varphi(t)) \varphi'(t) dt$$
\end{theorem}

% ================================
% 第二章：线性代数
% ================================

\section{线性代数}

\subsection{向量空间}

\begin{definition}[向量空间]
设 $V$ 是一个非空集合，$F$ 是一个域（通常取 $\mathbb{R}$ 或 $\mathbb{C}$）。若 $V$ 上定义了加法运算 $+: V \times V \to V$ 和标量乘法运算 $\cdot: F \times V \to V$，满足以下公理，则称 $(V, F, +, \cdot)$ 是 $F$ 上的向量空间：

\textbf{加法公理：}
\begin{enumerate}
\item 交换律：$\mathbf{u} + \mathbf{v} = \mathbf{v} + \mathbf{u}$
\item 结合律：$(\mathbf{u} + \mathbf{v}) + \mathbf{w} = \mathbf{u} + (\mathbf{v} + \mathbf{w})$
\item 零元素：存在 $\mathbf{0} \in V$ 使得 $\mathbf{v} + \mathbf{0} = \mathbf{v}$
\item 逆元素：对于每个 $\mathbf{v} \in V$，存在 $-\mathbf{v} \in V$ 使得 $\mathbf{v} + (-\mathbf{v}) = \mathbf{0}$
\end{enumerate}

\textbf{标量乘法公理：}
\begin{enumerate}
\item $a(b\mathbf{v}) = (ab)\mathbf{v}$
\item $1\mathbf{v} = \mathbf{v}$
\item $(a + b)\mathbf{v} = a\mathbf{v} + b\mathbf{v}$
\item $a(\mathbf{u} + \mathbf{v}) = a\mathbf{u} + a\mathbf{v}$
\end{enumerate}
\end{definition}

\begin{example}
常见的向量空间：
\begin{enumerate}
\item $\mathbb{R}^n$：$n$ 维实向量空间
\item $\mathbb{C}^n$：$n$ 维复向量空间
\item $M_{m \times n}(\mathbb{R})$：$m \times n$ 实矩阵空间
\item $P_n(\mathbb{R})$：次数不超过 $n$ 的实多项式空间
\item $C[a, b]$：$[a, b]$ 上连续函数空间
\end{enumerate}
\end{example}

\begin{definition}[线性组合与生成空间]
设 $V$ 是向量空间，$S = \{\mathbf{v}_1, \mathbf{v}_2, \ldots, \mathbf{v}_k\} \subseteq V$。

向量 $\mathbf{v}$ 是 $S$ 的线性组合，若存在标量 $c_1, c_2, \ldots, c_k$ 使得：
$$\mathbf{v} = c_1\mathbf{v}_1 + c_2\mathbf{v}_2 + \cdots + c_k\mathbf{v}_k$$

$S$ 生成的子空间（或张成空间）定义为：
$$\text{span}(S) = \{c_1\mathbf{v}_1 + c_2\mathbf{v}_2 + \cdots + c_k\mathbf{v}_k : c_i \in F\}$$
\end{definition}

\begin{definition}[线性无关与线性相关]
向量组 $\{\mathbf{v}_1, \mathbf{v}_2, \ldots, \mathbf{v}_k\}$ 称为线性无关的，若方程：
$$c_1\mathbf{v}_1 + c_2\mathbf{v}_2 + \cdots + c_k\mathbf{v}_k = \mathbf{0}$$
仅有平凡解 $c_1 = c_2 = \cdots = c_k = 0$。

否则称为线性相关的。
\end{definition}

\begin{definition}[基与维数]
向量空间 $V$ 的基是一个线性无关的向量组 $\mathcal{B} = \{\mathbf{v}_1, \mathbf{v}_2, \ldots, \mathbf{v}_n\}$，满足 $V = \text{span}(\mathcal{B})$。

基中向量的个数称为 $V$ 的维数，记为 $\dim V$。
\end{definition}

\begin{theorem}[基的存在性与唯一性]
\begin{enumerate}
\item 每个有限维向量空间都有基
\item 向量空间的任意两个基含有相同数目的向量
\end{enumerate}
\end{theorem}

\begin{theorem}[维数公式]
设 $W_1, W_2$ 是有限维向量空间 $V$ 的子空间，则：
$$\dim(W_1 + W_2) = \dim W_1 + \dim W_2 - \dim(W_1 \cap W_2)$$
\end{theorem}

\subsection{线性映射}

\begin{definition}[线性映射]
设 $V, W$ 是域 $F$ 上的向量空间。映射 $T: V \to W$ 称为线性映射（或线性变换），若满足：
\begin{enumerate}
\item $T(\mathbf{u} + \mathbf{v}) = T(\mathbf{u}) + T(\mathbf{v})$
\item $T(c\mathbf{v}) = cT(\mathbf{v})$
\end{enumerate}
对所有 $\mathbf{u}, \mathbf{v} \in V$ 和 $c \in F$ 成立。
\end{definition}

\begin{definition}[核与像]
设 $T: V \to W$ 是线性映射。

$T$ 的核（零空间）：$\ker T = \{\mathbf{v} \in V : T(\mathbf{v}) = \mathbf{0}\}$

$T$ 的像（值域）：$\text{Im } T = \{T(\mathbf{v}) : \mathbf{v} \in V\}$
\end{definition}

\begin{theorem}[秩-零化度定理]
设 $T: V \to W$ 是有限维向量空间之间的线性映射，则：
$$\dim V = \dim(\ker T) + \dim(\text{Im } T)$$
即 $\dim V = \text{nullity}(T) + \text{rank}(T)$。
\end{theorem}

\begin{proof}
设 $\dim(\ker T) = k$，取 $\ker T$ 的基 $\{\mathbf{u}_1, \ldots, \mathbf{u}_k\}$。

将其扩展为 $V$ 的基 $\{\mathbf{u}_1, \ldots, \mathbf{u}_k, \mathbf{v}_1, \ldots, \mathbf{v}_r\}$，其中 $k + r = \dim V$。

我们证明 $\{T(\mathbf{v}_1), \ldots, T(\mathbf{v}_r)\}$ 是 $\text{Im } T$ 的基。

\textbf{生成性：} 对于任意 $T(\mathbf{w}) \in \text{Im } T$，设 $\mathbf{w} = \sum_{i=1}^{k} a_i\mathbf{u}_i + \sum_{j=1}^{r} b_j\mathbf{v}_j$，则：
$$T(\mathbf{w}) = \sum_{j=1}^{r} b_j T(\mathbf{v}_j)$$

\textbf{线性无关性：} 设 $\sum_{j=1}^{r} c_j T(\mathbf{v}_j) = \mathbf{0}$，则 $T(\sum_{j=1}^{r} c_j\mathbf{v}_j) = \mathbf{0}$。

因此 $\sum_{j=1}^{r} c_j\mathbf{v}_j \in \ker T$，可表为 $\sum_{i=1}^{k} a_i\mathbf{u}_i$。

由 $\{\mathbf{u}_1, \ldots, \mathbf{u}_k, \mathbf{v}_1, \ldots, \mathbf{v}_r\}$ 线性无关，$c_1 = \cdots = c_r = 0$。

因此 $\dim(\text{Im } T) = r$，$\dim V = k + r = \dim(\ker T) + \dim(\text{Im } T)$。
\end{proof}

\subsection{矩阵理论}

\begin{definition}[矩阵]
$m \times n$ 矩阵是一个 $m$ 行 $n$ 列的数表：
$$A = \begin{pmatrix}
a_{11} & a_{12} & \cdots & a_{1n} \\
a_{21} & a_{22} & \cdots & a_{2n} \\
\vdots & \vdots & \ddots & \vdots \\
a_{m1} & a_{m2} & \cdots & a_{mn}
\end{pmatrix}$$
简记为 $A = (a_{ij})_{m \times n}$ 或 $A \in M_{m \times n}(F)$。
\end{definition}

\begin{definition}[矩阵运算]
设 $A = (a_{ij})_{m \times n}$，$B = (b_{ij})_{m \times n}$，$C = (c_{jk})_{n \times p}$，$\lambda \in F$。

\textbf{加法：} $(A + B)_{ij} = a_{ij} + b_{ij}$

\textbf{标量乘法：} $(\lambda A)_{ij} = \lambda a_{ij}$

\textbf{矩阵乘法：} $(AC)_{ik} = \sum_{j=1}^{n} a_{ij} c_{jk}$

\textbf{转置：} $(A^T)_{ij} = a_{ji}$
\end{definition}

\begin{theorem}[矩阵乘法性质]
\begin{enumerate}
\item 结合律：$(AB)C = A(BC)$
\item 分配律：$A(B + C) = AB + AC$，$(A + B)C = AC + BC$
\item $(AB)^T = B^T A^T$
\item 一般不满足交换律：$AB \neq BA$
\end{enumerate}
\end{theorem}

\begin{definition}[行列式]
$n$ 阶方阵 $A$ 的行列式定义为：
$$\det A = \sum_{\sigma \in S_n} \text{sgn}(\sigma) \prod_{i=1}^{n} a_{i,\sigma(i)}$$
其中 $S_n$ 是 $\{1, 2, \ldots, n\}$ 的所有排列组成的对称群，$\text{sgn}(\sigma)$ 是排列的符号。
\end{definition}

\begin{theorem}[行列式性质]
\begin{enumerate}
\item $\det(A^T) = \det A$
\item $\det(AB) = \det A \cdot \det B$
\item $\det(cA) = c^n \det A$（$A$ 是 $n$ 阶矩阵）
\item 行列式关于每行（列）是线性的
\item 交换两行（列），行列式变号
\item 两行（列）相同，行列式为零
\item $A$ 可逆当且仅当 $\det A \neq 0$
\end{enumerate}
\end{theorem}

\begin{definition}[逆矩阵]
方阵 $A$ 称为可逆的（或非奇异的），若存在方阵 $B$ 使得 $AB = BA = I$。这样的 $B$ 称为 $A$ 的逆矩阵，记为 $A^{-1}$。
\end{definition}

\begin{theorem}[逆矩阵公式]
若 $\det A \neq 0$，则：
$$A^{-1} = \frac{1}{\det A} \text{adj}(A)$$
其中 $\text{adj}(A)$ 是 $A$ 的伴随矩阵，$(\text{adj}(A))_{ij} = (-1)^{i+j} M_{ji}$，$M_{ji}$ 是 $a_{ji}$ 的余子式。
\end{theorem}

\begin{theorem}[克拉默法则]
设线性方程组 $A\mathbf{x} = \mathbf{b}$，其中 $A$ 是 $n$ 阶可逆矩阵。则方程组有唯一解：
$$x_i = \frac{\det A_i}{\det A}$$
其中 $A_i$ 是将 $A$ 的第 $i$ 列替换为 $\mathbf{b}$ 得到的矩阵。
\end{theorem}

\subsection{特征值与特征向量}

\begin{definition}[特征值与特征向量]
设 $A$ 是 $n$ 阶方阵。若存在标量 $\lambda$ 和非零向量 $\mathbf{v}$ 使得：
$$A\mathbf{v} = \lambda\mathbf{v}$$
则称 $\lambda$ 是 $A$ 的特征值，$\mathbf{v}$ 是对应于 $\lambda$ 的特征向量。
\end{definition}

\begin{definition}[特征多项式]
矩阵 $A$ 的特征多项式定义为：
$$p_A(\lambda) = \det(\lambda I - A)$$
特征值就是特征多项式的根。
\end{definition}

\begin{theorem}[Cayley-Hamilton 定理]
设 $p_A(\lambda)$ 是矩阵 $A$ 的特征多项式，则 $p_A(A) = O$（零矩阵）。
\end{theorem}

\begin{theorem}[对角化条件]
$n$ 阶矩阵 $A$ 可对角化的充要条件是 $A$ 有 $n$ 个线性无关的特征向量。

等价条件：对于每个特征值 $\lambda$，几何重数等于代数重数。
\end{theorem}

\begin{theorem}[实对称矩阵的对角化]
实对称矩阵 $A$ 一定可以正交对角化，即存在正交矩阵 $P$（$P^T P = I$）使得：
$$P^T A P = D$$
其中 $D$ 是对角矩阵，对角元素是 $A$ 的特征值。
\end{theorem}

\begin{theorem}[谱定理]
设 $A$ 是 $n$ 阶实对称矩阵，$\lambda_1, \ldots, \lambda_n$ 是其特征值（含重复），$\mathbf{v}_1, \ldots, \mathbf{v}_n$ 是对应的标准正交特征向量。则：
$$A = \sum_{i=1}^{n} \lambda_i \mathbf{v}_i \mathbf{v}_i^T$$
这称为 $A$ 的谱分解。
\end{theorem}

% ================================
% 第三章：概率论与统计
% ================================

\section{概率论与数理统计}

\subsection{概率空间}

\begin{definition}[概率空间]
概率空间是三元组 $(\Omega, \mathcal{F}, P)$，其中：
\begin{enumerate}
\item $\Omega$ 是样本空间（所有可能结果的集合）
\item $\mathcal{F}$ 是 $\Omega$ 上的 $\sigma$-代数（事件空间）
\item $P: \mathcal{F} \to [0, 1]$ 是概率测度，满足：
    \begin{enumerate}
    \item $P(\Omega) = 1$
    \item 可列可加性：对于互不相交的事件序列 $\{A_n\}$，$P(\bigcup_{n=1}^{\infty} A_n) = \sum_{n=1}^{\infty} P(A_n)$
    \end{enumerate}
\end{enumerate}
\end{definition}

\begin{theorem}[概率的基本性质]
\begin{enumerate}
\item $P(\emptyset) = 0$
\item $P(A^c) = 1 - P(A)$
\item $A \subseteq B \Rightarrow P(A) \leq P(B)$
\item $P(A \cup B) = P(A) + P(B) - P(A \cap B)$
\item 次可加性：$P(\bigcup_{i=1}^{n} A_i) \leq \sum_{i=1}^{n} P(A_i)$
\end{enumerate}
\end{theorem}

\begin{definition}[条件概率]
设 $P(B) > 0$，事件 $A$ 在事件 $B$ 发生条件下的条件概率定义为：
$$P(A|B) = \frac{P(A \cap B)}{P(B)}$$
\end{definition}

\begin{theorem}[乘法公式]
$$P(A \cap B) = P(A|B) \cdot P(B) = P(B|A) \cdot P(A)$$

更一般地：
$$P(A_1 \cap A_2 \cap \cdots \cap A_n) = P(A_1) \cdot P(A_2|A_1) \cdot P(A_3|A_1 \cap A_2) \cdots P(A_n|A_1 \cap \cdots \cap A_{n-1})$$
\end{theorem}

\begin{theorem}[全概率公式]
设 $\{B_1, B_2, \ldots, B_n\}$ 是样本空间 $\Omega$ 的一个划分，$P(B_i) > 0$，则：
$$P(A) = \sum_{i=1}^{n} P(A|B_i) P(B_i)$$
\end{theorem}

\begin{theorem}[贝叶斯公式]
设 $\{B_1, B_2, \ldots, B_n\}$ 是样本空间 $\Omega$ 的一个划分，$P(B_i) > 0$，$P(A) > 0$，则：
$$P(B_j|A) = \frac{P(A|B_j) P(B_j)}{\sum_{i=1}^{n} P(A|B_i) P(B_i)}$$
\end{theorem}

\begin{definition}[事件独立性]
事件 $A$ 和 $B$ 称为独立的，若 $P(A \cap B) = P(A) \cdot P(B)$。

事件族 $\{A_i : i \in I\}$ 称为相互独立的，若对于任意有限子集 $J \subseteq I$：
$$P\left(\bigcap_{j \in J} A_j\right) = \prod_{j \in J} P(A_j)$$
\end{definition}

\subsection{随机变量}

\begin{definition}[随机变量]
设 $(\Omega, \mathcal{F}, P)$ 是概率空间。随机变量是可测映射 $X: \Omega \to \mathbb{R}$，即对于任意 Borel 集 $B$，$X^{-1}(B) \in \mathcal{F}$。
\end{definition}

\begin{definition}[分布函数]
随机变量 $X$ 的分布函数定义为：
$$F_X(x) = P(X \leq x), \quad x \in \mathbb{R}$$
\end{definition}

\begin{theorem}[分布函数的性质]
\begin{enumerate}
\item 单调不减：$x_1 < x_2 \Rightarrow F(x_1) \leq F(x_2)$
\item 右连续：$\lim_{x \to a^+} F(x) = F(a)$
\item $\lim_{x \to -\infty} F(x) = 0$，$\lim_{x \to +\infty} F(x) = 1$
\end{enumerate}
\end{theorem}

\begin{definition}[离散随机变量]
若随机变量 $X$ 只取有限个或可数个值 $x_1, x_2, \ldots$，则称 $X$ 为离散随机变量。其概率质量函数定义为：
$$p_X(x_i) = P(X = x_i)$$
\end{definition}

\begin{definition}[连续随机变量]
若存在非负函数 $f_X(x)$ 使得：
$$F_X(x) = \int_{-\infty}^{x} f_X(t) dt$$
则称 $X$ 为连续随机变量，$f_X(x)$ 称为 $X$ 的概率密度函数。
\end{definition}

\textbf{常见离散分布：}

\begin{enumerate}
\item \textbf{伯努利分布} $B(1, p)$：$P(X = 1) = p$，$P(X = 0) = 1 - p$

\item \textbf{二项分布} $B(n, p)$：
$$P(X = k) = \binom{n}{k} p^k (1-p)^{n-k}, \quad k = 0, 1, \ldots, n$$
$E[X] = np$，$\text{Var}(X) = np(1-p)$

\item \textbf{泊松分布} $\text{Poisson}(\lambda)$：
$$P(X = k) = \frac{\lambda^k e^{-\lambda}}{k!}, \quad k = 0, 1, 2, \ldots$$
$E[X] = \lambda$，$\text{Var}(X) = \lambda$

\item \textbf{几何分布} $\text{Geo}(p)$：
$$P(X = k) = (1-p)^{k-1} p, \quad k = 1, 2, \ldots$$
$E[X] = 1/p$，$\text{Var}(X) = (1-p)/p^2$
\end{enumerate}

\textbf{常见连续分布：}

\begin{enumerate}
\item \textbf{均匀分布} $U(a, b)$：
$$f(x) = \begin{cases} \frac{1}{b-a} & a \leq x \leq b \\ 0 & \text{其他} \end{cases}$$
$E[X] = \frac{a+b}{2}$，$\text{Var}(X) = \frac{(b-a)^2}{12}$

\item \textbf{指数分布} $\text{Exp}(\lambda)$：
$$f(x) = \begin{cases} \lambda e^{-\lambda x} & x \geq 0 \\ 0 & x < 0 \end{cases}$$
$E[X] = 1/\lambda$，$\text{Var}(X) = 1/\lambda^2$

\item \textbf{正态分布} $N(\mu, \sigma^2)$：
$$f(x) = \frac{1}{\sqrt{2\pi}\sigma} \exp\left(-\frac{(x-\mu)^2}{2\sigma^2}\right)$$
$E[X] = \mu$，$\text{Var}(X) = \sigma^2$

\item \textbf{Gamma 分布} $\Gamma(\alpha, \beta)$：
$$f(x) = \frac{\beta^\alpha}{\Gamma(\alpha)} x^{\alpha-1} e^{-\beta x}, \quad x > 0$$
$E[X] = \alpha/\beta$，$\text{Var}(X) = \alpha/\beta^2$
\end{enumerate}

\subsection{数学期望与方差}

\begin{definition}[数学期望]
\textbf{离散情形：}
$$E[X] = \sum_{i} x_i P(X = x_i)$$

\textbf{连续情形：}
$$E[X] = \int_{-\infty}^{\infty} x f_X(x) dx$$
\end{definition}

\begin{theorem}[期望的性质]
\begin{enumerate}
\item 线性性：$E[aX + bY] = aE[X] + bE[Y]$
\item 若 $X \geq 0$，则 $E[X] \geq 0$
\item 若 $X, Y$ 独立，则 $E[XY] = E[X] \cdot E[Y]$
\end{enumerate}
\end{theorem}

\begin{definition}[方差与标准差]
$$\text{Var}(X) = E[(X - E[X])^2] = E[X^2] - (E[X])^2$$
标准差：$\sigma_X = \sqrt{\text{Var}(X)}$
\end{definition}

\begin{theorem}[方差的性质]
\begin{enumerate}
\item $\text{Var}(aX + b) = a^2 \text{Var}(X)$
\item 若 $X, Y$ 独立，则 $\text{Var}(X + Y) = \text{Var}(X) + \text{Var}(Y)$
\item $\text{Var}(X) = 0$ 当且仅当 $P(X = E[X]) = 1$
\end{enumerate}
\end{theorem}

\begin{definition}[协方差与相关系数]
$$\text{Cov}(X, Y) = E[(X - E[X])(Y - E[Y])] = E[XY] - E[X]E[Y]$$

$$\rho_{XY} = \frac{\text{Cov}(X, Y)}{\sigma_X \sigma_Y}$$
\end{definition}

\begin{theorem}[协方差的性质]
\begin{enumerate}
\item $\text{Cov}(X, X) = \text{Var}(X)$
\item $\text{Cov}(X, Y) = \text{Cov}(Y, X)$
\item $\text{Cov}(aX, bY) = ab \cdot \text{Cov}(X, Y)$
\item $\text{Var}(X + Y) = \text{Var}(X) + \text{Var}(Y) + 2\text{Cov}(X, Y)$
\item $|\rho_{XY}| \leq 1$，$|\rho_{XY}| = 1$ 当且仅当 $Y = aX + b$（a.s.）
\end{enumerate}
\end{theorem}

\subsection{大数定律与中心极限定理}

\begin{theorem}[马尔可夫不等式]
设 $X$ 是非负随机变量，$a > 0$，则：
$$P(X \geq a) \leq \frac{E[X]}{a}$$
\end{theorem}

\begin{theorem}[切比雪夫不等式]
设 $X$ 的期望和方差存在，则对于任意 $\epsilon > 0$：
$$P(|X - E[X]| \geq \epsilon) \leq \frac{\text{Var}(X)}{\epsilon^2}$$
\end{theorem}

\begin{theorem}[弱大数定律]
设 $X_1, X_2, \ldots$ 是独立同分布的随机变量序列，$E[X_i] = \mu$，$\text{Var}(X_i) = \sigma^2 < \infty$。令 $\bar{X}_n = \frac{1}{n}\sum_{i=1}^{n} X_i$，则对于任意 $\epsilon > 0$：
$$\lim_{n \to \infty} P(|\bar{X}_n - \mu| \geq \epsilon) = 0$$
\end{theorem}

\begin{theorem}[强大数定律]
设 $X_1, X_2, \ldots$ 是独立同分布的随机变量序列，$E[|X_i|] < \infty$，$E[X_i] = \mu$。则：
$$P\left(\lim_{n \to \infty} \bar{X}_n = \mu\right) = 1$$
\end{theorem}

\begin{theorem}[中心极限定理 (Lindeberg-Lévy)]
设 $X_1, X_2, \ldots$ 是独立同分布的随机变量序列，$E[X_i] = \mu$，$\text{Var}(X_i) = \sigma^2 \in (0, \infty)$。令 $S_n = \sum_{i=1}^{n} X_i$，则：
$$\frac{S_n - n\mu}{\sigma\sqrt{n}} \xrightarrow{d} N(0, 1)$$
即对于任意 $x \in \mathbb{R}$：
$$\lim_{n \to \infty} P\left(\frac{S_n - n\mu}{\sigma\sqrt{n}} \leq x\right) = \Phi(x) = \frac{1}{\sqrt{2\pi}} \int_{-\infty}^{x} e^{-t^2/2} dt$$
\end{theorem}

% ================================
% 第四章：复分析
% ================================

\section{复分析}

\subsection{复数与复平面}

\begin{definition}[复数]
复数是形如 $z = x + iy$ 的数，其中 $x, y \in \mathbb{R}$，$i$ 是虚数单位，满足 $i^2 = -1$。

$x$ 称为 $z$ 的实部，记为 $\text{Re}(z)$；$y$ 称为 $z$ 的虚部，记为 $\text{Im}(z)$。

复数的模（绝对值）定义为：$|z| = \sqrt{x^2 + y^2}$

复数的辐角定义为：$\arg(z) = \arctan(y/x)$（需考虑象限）

复数的共轭定义为：$\bar{z} = x - iy$
\end{definition}

\begin{theorem}[复数的极坐标表示]
任意非零复数 $z$ 可表示为：
$$z = r(\cos\theta + i\sin\theta) = re^{i\theta}$$
其中 $r = |z|$，$\theta = \arg(z)$。

欧拉公式：$e^{i\theta} = \cos\theta + i\sin\theta$
\end{theorem}

\begin{theorem}[复数运算性质]
设 $z_1 = x_1 + iy_1$，$z_2 = x_2 + iy_2$，则：
\begin{align}
z_1 + z_2 &= (x_1 + x_2) + i(y_1 + y_2) \\
z_1 \cdot z_2 &= (x_1x_2 - y_1y_2) + i(x_1y_2 + x_2y_1) \\
|z_1 z_2| &= |z_1| \cdot |z_2| \\
\arg(z_1 z_2) &= \arg(z_1) + \arg(z_2) \\
|z_1 + z_2| &\leq |z_1| + |z_2| \quad \text{（三角不等式）}
\end{align}
\end{theorem}

\subsection{复变函数}

\begin{definition}[复变函数]
设 $D \subseteq \mathbb{C}$ 是复平面的子集。复变函数 $f: D \to \mathbb{C}$ 将每个 $z \in D$ 映射到 $w = f(z) \in \mathbb{C}$。

将 $z = x + iy$，$w = u + iv$，则 $f$ 可表示为：
$$f(z) = u(x, y) + iv(x, y)$$
其中 $u, v: \mathbb{R}^2 \to \mathbb{R}$ 是实函数。
\end{definition}

\begin{definition}[复函数的极限]
设 $f$ 在 $z_0$ 的去心邻域内有定义。若对于任意 $\epsilon > 0$，存在 $\delta > 0$，使得当 $0 < |z - z_0| < \delta$ 时，有 $|f(z) - L| < \epsilon$，则称 $L$ 为 $f$ 在 $z_0$ 处的极限，记为：
$$\lim_{z \to z_0} f(z) = L$$
\end{definition}

\begin{definition}[复函数的连续性]
若 $\lim_{z \to z_0} f(z) = f(z_0)$，则称 $f$ 在 $z_0$ 处连续。
\end{definition}

\begin{definition}[复函数的导数（解析性）]
若极限
$$f'(z_0) = \lim_{h \to 0} \frac{f(z_0 + h) - f(z_0)}{h}$$
存在，则称 $f$ 在 $z_0$ 处可导（或复可微）。

若 $f$ 在区域 $D$ 内每一点都可导，则称 $f$ 在 $D$ 内解析（或全纯）。
\end{definition}

\begin{theorem}[柯西-黎曼方程]
设 $f(z) = u(x, y) + iv(x, y)$。$f$ 在 $z_0 = x_0 + iy_0$ 处可导的充要条件是：
\begin{enumerate}
\item $u, v$ 在 $(x_0, y_0)$ 处可微
\item 柯西-黎曼方程成立：
$$\frac{\partial u}{\partial x} = \frac{\partial v}{\partial y}, \quad \frac{\partial u}{\partial y} = -\frac{\partial v}{\partial x}$$
\end{enumerate}

此时：
$$f'(z_0) = \frac{\partial u}{\partial x} + i\frac{\partial v}{\partial x} = \frac{\partial v}{\partial y} - i\frac{\partial u}{\partial y}$$
\end{theorem}

\begin{theorem}[解析函数的性质]
设 $f$ 在区域 $D$ 内解析，则：
\begin{enumerate}
\item $f$ 在 $D$ 内无穷次可微
\item $f$ 在 $D$ 内可展开为幂级数
\item $u = \text{Re}(f)$ 和 $v = \text{Im}(f)$ 都是调和函数，即满足拉普拉斯方程：
$$\Delta u = \frac{\partial^2 u}{\partial x^2} + \frac{\partial^2 u}{\partial y^2} = 0$$
\end{enumerate}
\end{theorem}

\subsection{复积分}

\begin{definition}[复积分]
设 $\gamma: [a, b] \to \mathbb{C}$ 是光滑曲线，$f$ 在 $\gamma$ 上连续。$f$ 沿 $\gamma$ 的积分定义为：
$$\int_\gamma f(z) dz = \int_a^b f(\gamma(t)) \gamma'(t) dt$$
\end{definition}

\begin{theorem}[柯西积分定理]
设 $f$ 在单连通区域 $D$ 内解析，$\gamma$ 是 $D$ 内的闭曲线，则：
$$\oint_\gamma f(z) dz = 0$$
\end{theorem}

\begin{theorem}[柯西积分公式]
设 $f$ 在单连通区域 $D$ 内解析，$\gamma$ 是 $D$ 内包围点 $z_0$ 的简单闭曲线，则：
$$f(z_0) = \frac{1}{2\pi i} \oint_\gamma \frac{f(z)}{z - z_0} dz$$

更一般地，$f$ 的 $n$ 阶导数为：
$$f^{(n)}(z_0) = \frac{n!}{2\pi i} \oint_\gamma \frac{f(z)}{(z - z_0)^{n+1}} dz$$
\end{theorem}

\begin{theorem}[刘维尔定理]
若 $f$ 是整函数（在整个复平面上解析）且有界，则 $f$ 是常函数。
\end{theorem}

\begin{theorem}[代数基本定理]
每个次数 $n \geq 1$ 的复系数多项式在 $\mathbb{C}$ 中恰有 $n$ 个根（计重数）。
\end{theorem}

\subsection{级数与留数}

\begin{definition}[幂级数]
形如 $\sum_{n=0}^{\infty} a_n (z - z_0)^n$ 的级数称为以 $z_0$ 为中心的幂级数。

收敛半径定义为：
$$R = \frac{1}{\limsup_{n \to \infty} |a_n|^{1/n}}$$

级数在 $|z - z_0| < R$ 内绝对收敛，在 $|z - z_0| > R$ 内发散。
\end{definition}

\begin{theorem}[泰勒级数]
若 $f$ 在区域 $D$ 内解析，$z_0 \in D$，$B_R(z_0) \subseteq D$，则在 $B_R(z_0)$ 内：
$$f(z) = \sum_{n=0}^{\infty} \frac{f^{(n)}(z_0)}{n!} (z - z_0)^n$$
\end{theorem}

\begin{definition}[洛朗级数]
若 $f$ 在圆环 $r < |z - z_0| < R$ 内解析，则 $f$ 可展开为洛朗级数：
$$f(z) = \sum_{n=-\infty}^{\infty} a_n (z - z_0)^n$$
其中：
$$a_n = \frac{1}{2\pi i} \oint_\gamma \frac{f(z)}{(z - z_0)^{n+1}} dz$$
$\gamma$ 是圆环内任意一条绕 $z_0$ 的简单闭曲线。
\end{definition}

\begin{definition}[奇点分类]
设 $z_0$ 是 $f$ 的孤立奇点，$f$ 在 $z_0$ 处的洛朗展开为 $\sum_{n=-\infty}^{\infty} a_n (z - z_0)^n$。

\begin{enumerate}
\item 可去奇点：$a_n = 0$ 对所有 $n < 0$ 成立
\item $m$ 阶极点：$a_{-m} \neq 0$，$a_n = 0$ 对所有 $n < -m$ 成立
\item 本质奇点：有无穷多个 $a_n$（$n < 0$）非零
\end{enumerate}
\end{definition}

\begin{definition}[留数]
$f$ 在孤立奇点 $z_0$ 处的留数定义为洛朗级数中 $(z - z_0)^{-1}$ 的系数：
$$\text{Res}(f, z_0) = a_{-1} = \frac{1}{2\pi i} \oint_\gamma f(z) dz$$
\end{definition}

\begin{theorem}[留数定理]
设 $f$ 在区域 $D$ 内除有限个孤立奇点 $z_1, \ldots, z_n$ 外解析，$\gamma$ 是 $D$ 内包围所有奇点的简单闭曲线，则：
$$\oint_\gamma f(z) dz = 2\pi i \sum_{k=1}^{n} \text{Res}(f, z_k)$$
\end{theorem}

\begin{theorem}[留数计算公式]
设 $z_0$ 是 $f$ 的 $m$ 阶极点，则：
$$\text{Res}(f, z_0) = \frac{1}{(m-1)!} \lim_{z \to z_0} \frac{d^{m-1}}{dz^{m-1}} [(z - z_0)^m f(z)]$$

特别地，若 $z_0$ 是单极点：
$$\text{Res}(f, z_0) = \lim_{z \to z_0} (z - z_0) f(z)$$
\end{theorem}

% ================================
% 第五章：拓扑学基础
% ================================

\section{拓扑学基础}

\subsection{拓扑空间}

\begin{definition}[拓扑空间]
设 $X$ 是一个集合，$\mathcal{T} \subseteq \mathcal{P}(X)$（$X$ 的幂集的子集）。若 $\mathcal{T}$ 满足：
\begin{enumerate}
\item $\emptyset, X \in \mathcal{T}$
\item $\mathcal{T}$ 对任意并封闭：$\{U_\alpha\}_{\alpha \in I} \subseteq \mathcal{T} \Rightarrow \bigcup_{\alpha \in I} U_\alpha \in \mathcal{T}$
\item $\mathcal{T}$ 对有限交封闭：$U_1, U_2 \in \mathcal{T} \Rightarrow U_1 \cap U_2 \in \mathcal{T}$
\end{enumerate}
则称 $(X, \mathcal{T})$ 是一个拓扑空间，$\mathcal{T}$ 称为 $X$ 上的拓扑，$\mathcal{T}$ 中的元素称为开集。
\end{definition}

\begin{example}
常见拓扑：
\begin{enumerate}
\item 离散拓扑：$\mathcal{T} = \mathcal{P}(X)$
\item 平凡拓扑：$\mathcal{T} = \{\emptyset, X\}$
\item 欧氏拓扑：$\mathbb{R}^n$ 上由开球生成的拓扑
\item 余有限拓扑：$U \in \mathcal{T}$ 当且仅当 $X \setminus U$ 有限或 $U = \emptyset$
\end{enumerate}
\end{example}

\begin{definition}[闭集、闭包、内部、边界]
\begin{enumerate}
\item 闭集：开集的补集
\item 闭包：$\bar{A} = \bigcap\{F : A \subseteq F, F \text{ 是闭集}\}$
\item 内部：$A^\circ = \bigcup\{U : U \subseteq A, U \text{ 是开集}\}$
\item 边界：$\partial A = \bar{A} \setminus A^\circ$
\end{enumerate}
\end{definition}

\begin{definition}[基与子基]
拓扑 $\mathcal{T}$ 的基是集族 $\mathcal{B} \subseteq \mathcal{T}$，满足每个 $U \in \mathcal{T}$ 都可表示为 $\mathcal{B}$ 中某些元素的并。

子基是集族 $\mathcal{S}$，其有限交的全体构成 $\mathcal{T}$ 的基。
\end{definition}

\begin{definition}[连续映射]
设 $(X, \mathcal{T}_X)$ 和 $(Y, \mathcal{T}_Y)$ 是拓扑空间。映射 $f: X \to Y$ 称为连续的，若对于每个 $V \in \mathcal{T}_Y$，$f^{-1}(V) \in \mathcal{T}_X$。
\end{definition}

\begin{definition}[同胚]
若 $f: X \to Y$ 是双射，且 $f$ 和 $f^{-1}$ 都连续，则称 $f$ 是同胚，$X$ 和 $Y$ 称为同胚的。
\end{definition}

\subsection{紧性与连通性}

\begin{definition}[紧空间]
拓扑空间 $X$ 称为紧的，若 $X$ 的每个开覆盖都有有限子覆盖。即：若 $\{U_\alpha\}_{\alpha \in I}$ 满足 $X = \bigcup_{\alpha \in I} U_\alpha$，则存在有限集 $J \subseteq I$ 使得 $X = \bigcup_{\alpha \in J} U_\alpha$。
\end{definition}

\begin{theorem}[海涅-博雷尔定理]
$\mathbb{R}^n$ 的子集 $K$ 紧当且仅当 $K$ 有界且闭。
\end{theorem}

\begin{theorem}[紧空间的性质]
\begin{enumerate}
\item 紧空间的闭子集是紧的
\item 紧空间到豪斯多夫空间的连续映射的像是紧的
\item 豪斯多夫空间的紧子集是闭的
\item 紧豪斯多夫空间到豪斯多夫空间的连续双射是同胚
\end{enumerate}
\end{theorem}

\begin{definition}[连通空间]
拓扑空间 $X$ 称为连通的，若 $X$ 不能表示为两个非空不相交开集的并。

等价地，$X$ 连通当且仅当 $X$ 中唯一既开又闭的集合是 $\emptyset$ 和 $X$。
\end{definition}

\begin{definition}[道路连通]
$X$ 称为道路连通的，若对于任意 $x, y \in X$，存在连续映射 $\gamma: [0, 1] \to X$ 使得 $\gamma(0) = x$，$\gamma(1) = y$。
\end{definition}

\begin{theorem}[连通性的性质]
\begin{enumerate}
\item 道路连通 $\Rightarrow$ 连通（反之不一定成立）
\item 连续映射保持连通性
\item $\mathbb{R}$ 的连通子集恰好是区间
\end{enumerate}
\end{theorem}

\subsection{分离公理}

\begin{definition}[分离公理]
\begin{enumerate}
\item $T_0$：对于任意两个不同点，至少有一个有不包含另一个的开邻域
\item $T_1$：对于任意两个不同点，每个都有不包含另一个的开邻域
\item $T_2$（豪斯多夫）：对于任意两个不同点，存在不相交的开邻域
\item $T_3$（正则）：$T_1$ + 点与不含它的闭集可用不相交开集分离
\item $T_4$（正规）：$T_1$ + 任意两个不相交闭集可用不相交开集分离
\end{enumerate}
\end{definition}

\begin{theorem}[Urysohn 引理]
设 $X$ 是正规空间，$A, B$ 是 $X$ 中不相交的闭集。则存在连续函数 $f: X \to [0, 1]$ 使得 $f|_A = 0$，$f|_B = 1$。
\end{theorem}

\begin{theorem}[Tietze 扩张定理]
设 $X$ 是正规空间，$A \subseteq X$ 是闭集，$f: A \to \mathbb{R}$ 连续。则存在连续函数 $F: X \to \mathbb{R}$ 使得 $F|_A = f$。
\end{theorem}

% ================================
% 第六章：微分方程
% ================================

\section{常微分方程}

\subsection{一阶微分方程}

\begin{definition}[一阶常微分方程]
一阶常微分方程是形如
$$F(x, y, y') = 0 \quad \text{或} \quad y' = f(x, y)$$
的方程，其中 $y = y(x)$ 是未知函数。
\end{definition}

\begin{theorem}[Picard-Lindelöf 存在唯一性定理]
设 $f(x, y)$ 在矩形区域 $R = \{(x, y) : |x - x_0| \leq a, |y - y_0| \leq b\}$ 上连续，且关于 $y$ 满足 Lipschitz 条件：
$$|f(x, y_1) - f(x, y_2)| \leq L|y_1 - y_2|$$
则初值问题
$$y' = f(x, y), \quad y(x_0) = y_0$$
在区间 $|x - x_0| \leq h$ 上有唯一解，其中 $h = \min(a, b/M)$，$M = \max_{R} |f|$。
\end{theorem}

\begin{example}[可分离变量方程]
形如 $\frac{dy}{dx} = g(x)h(y)$ 的方程，解法：
$$\int \frac{dy}{h(y)} = \int g(x) dx + C$$
\end{example}

\begin{example}[一阶线性方程]
形如 $y' + P(x)y = Q(x)$ 的方程，通解：
$$y = e^{-\int P dx}\left[\int Q e^{\int P dx} dx + C\right]$$
\end{example}

\begin{example}[伯努利方程]
形如 $y' + P(x)y = Q(x)y^n$（$n \neq 0, 1$）的方程。

令 $z = y^{1-n}$，则 $z' + (1-n)P(x)z = (1-n)Q(x)$，化为一阶线性方程。
\end{example}

\begin{example}[全微分方程]
形如 $M(x,y)dx + N(x,y)dy = 0$ 的方程，若 $\frac{\partial M}{\partial y} = \frac{\partial N}{\partial x}$，则存在势函数 $\Phi(x,y)$ 使得：
$$d\Phi = Mdx + Ndy$$
通解为 $\Phi(x, y) = C$。
\end{example}

\subsection{高阶线性微分方程}

\begin{definition}[$n$ 阶线性微分方程]
$$a_n(x)y^{(n)} + a_{n-1}(x)y^{(n-1)} + \cdots + a_1(x)y' + a_0(x)y = f(x)$$
若 $f(x) \equiv 0$，称为齐次方程；否则称为非齐次方程。
\end{definition}

\begin{theorem}[线性方程解的结构]
设 $L[y] = y^{(n)} + p_{n-1}(x)y^{(n-1)} + \cdots + p_0(x)y$。

\begin{enumerate}
\item 齐次方程 $L[y] = 0$ 的解构成 $n$ 维向量空间
\item 设 $y_1, \ldots, y_n$ 是 $L[y] = 0$ 的 $n$ 个线性无关解，则通解为：
$$y = C_1 y_1 + C_2 y_2 + \cdots + C_n y_n$$
\item 非齐次方程 $L[y] = f(x)$ 的通解为：
$$y = y_h + y_p$$
其中 $y_h$ 是齐次方程通解，$y_p$ 是非齐次方程的一个特解
\end{enumerate}
\end{theorem}

\begin{definition}[Wronskian 行列式]
设 $y_1, \ldots, y_n$ 是 $n$ 个 $n-1$ 次可微函数，其 Wronskian 定义为：
$$W(y_1, \ldots, y_n) = \begin{vmatrix}
y_1 & y_2 & \cdots & y_n \\
y_1' & y_2' & \cdots & y_n' \\
\vdots & \vdots & \ddots & \vdots \\
y_1^{(n-1)} & y_2^{(n-1)} & \cdots & y_n^{(n-1)}
\end{vmatrix}$$
\end{definition}

\begin{theorem}
齐次线性方程的 $n$ 个解 $y_1, \ldots, y_n$ 线性无关当且仅当 $W(y_1, \ldots, y_n) \not\equiv 0$。
\end{theorem}

\subsection{常系数线性微分方程}

\begin{theorem}[常系数齐次方程]
设方程为 $y^{(n)} + a_{n-1}y^{(n-1)} + \cdots + a_1y' + a_0y = 0$，特征方程为：
$$\lambda^n + a_{n-1}\lambda^{n-1} + \cdots + a_1\lambda + a_0 = 0$$

\begin{enumerate}
\item 若 $\lambda_0$ 是单实根，则对应解 $e^{\lambda_0 x}$
\item 若 $\lambda_0$ 是 $k$ 重实根，则对应解 $e^{\lambda_0 x}, xe^{\lambda_0 x}, \ldots, x^{k-1}e^{\lambda_0 x}$
\item 若 $\alpha \pm i\beta$ 是单复根，则对应解 $e^{\alpha x}\cos\beta x, e^{\alpha x}\sin\beta x$
\item 若 $\alpha \pm i\beta$ 是 $k$ 重复根，则对应解：
$$e^{\alpha x}\cos\beta x, xe^{\alpha x}\cos\beta x, \ldots, x^{k-1}e^{\alpha x}\cos\beta x$$
$$e^{\alpha x}\sin\beta x, xe^{\alpha x}\sin\beta x, \ldots, x^{k-1}e^{\alpha x}\sin\beta x$$
\end{enumerate}
\end{theorem}

\begin{theorem}[常数变易法]
设齐次方程 $L[y] = 0$ 的基本解组为 $y_1, \ldots, y_n$，则非齐次方程 $L[y] = f(x)$ 的特解为：
$$y_p = \sum_{i=1}^{n} y_i \int \frac{W_i}{W} f(x) dx$$
其中 $W = W(y_1, \ldots, y_n)$，$W_i$ 是将 $W$ 第 $i$ 列替换为 $(0, \ldots, 0, 1)^T$ 得到的行列式。
\end{theorem}

\subsection{微分方程组}

\begin{definition}[一阶线性微分方程组]
$$\mathbf{x}' = A(t)\mathbf{x} + \mathbf{f}(t)$$
其中 $\mathbf{x} = (x_1, \ldots, x_n)^T$，$A(t)$ 是 $n \times n$ 矩阵。
\end{definition}

\begin{theorem}[矩阵指数]
对于常系数齐次系统 $\mathbf{x}' = A\mathbf{x}$，解为：
$$\mathbf{x}(t) = e^{At}\mathbf{x}(0)$$
其中矩阵指数定义为：
$$e^{At} = \sum_{k=0}^{\infty} \frac{(At)^k}{k!} = I + At + \frac{A^2t^2}{2!} + \cdots$$
\end{theorem}

\begin{theorem}[常系数系统的解]
设 $A$ 的特征值为 $\lambda_1, \ldots, \lambda_n$，对应特征向量为 $\mathbf{v}_1, \ldots, \mathbf{v}_n$。

若 $A$ 可对角化，则通解为：
$$\mathbf{x}(t) = C_1 e^{\lambda_1 t}\mathbf{v}_1 + C_2 e^{\lambda_2 t}\mathbf{v}_2 + \cdots + C_n e^{\lambda_n t}\mathbf{v}_n$$
\end{theorem}

% ================================
% 第七章：数值分析
% ================================

\section{数值分析}

\subsection{误差分析}

\begin{definition}[误差]
设 $x^*$ 是精确值，$x$ 是近似值。

绝对误差：$E = x - x^*$

相对误差：$E_r = \frac{x - x^*}{x^*}$（$x^* \neq 0$）
\end{definition}

\begin{definition}[有效数字]
若 $x$ 是 $x^*$ 的近似值，且 $|x - x^*| \leq \frac{1}{2} \times 10^{-n}$，则称 $x$ 有 $n$ 位有效数字。
\end{definition}

\begin{theorem}[误差传播]
设 $y = f(x_1, \ldots, x_n)$，则：
$$\Delta y \approx \sum_{i=1}^{n} \frac{\partial f}{\partial x_i} \Delta x_i$$

$$|E_y| \lesssim \sum_{i=1}^{n} \left|\frac{\partial f}{\partial x_i}\right| |E_{x_i}|$$
\end{theorem}

\subsection{非线性方程求根}

\begin{theorem}[二分法]
设 $f$ 在 $[a, b]$ 上连续，$f(a)f(b) < 0$。二分法迭代：
\begin{enumerate}
\item 计算 $c = (a + b)/2$
\item 若 $f(c) = 0$，则 $c$ 是根
\item 若 $f(a)f(c) < 0$，取 $b = c$；否则取 $a = c$
\item 重复直到 $b - a < \epsilon$
\end{enumerate}

收敛速度：$|x_n - x^*| \leq \frac{b - a}{2^n}$
\end{theorem}

\begin{theorem}[牛顿法]
设 $f$ 可微。牛顿迭代：
$$x_{n+1} = x_n - \frac{f(x_n)}{f'(x_n)}$$

若 $x^*$ 是单根，$f''$ 在 $x^*$ 附近连续，且初值足够接近 $x^*$，则牛顿法二阶收敛：
$$|x_{n+1} - x^*| \leq C|x_n - x^*|^2$$
\end{theorem}

\begin{theorem}[割线法]
$$x_{n+1} = x_n - f(x_n) \frac{x_n - x_{n-1}}{f(x_n) - f(x_{n-1})}$$

收敛阶约为 $1.618$（黄金分割比）。
\end{theorem}

\begin{theorem}[不动点迭代]
设 $g: [a, b] \to [a, b]$，$|g'(x)| \leq L < 1$ 对所有 $x \in [a, b]$ 成立。则迭代：
$$x_{n+1} = g(x_n)$$
收敛到唯一不动点 $x^* = g(x^*)$，且：
$$|x_n - x^*| \leq \frac{L^n}{1 - L}|x_1 - x_0|$$
\end{theorem}

\subsection{插值与逼近}

\begin{definition}[Lagrange 插值]
给定 $n+1$ 个互异节点 $(x_0, y_0), \ldots, (x_n, y_n)$，Lagrange 插值多项式为：
$$L_n(x) = \sum_{i=0}^{n} y_i \prod_{j \neq i} \frac{x - x_j}{x_i - x_j}$$
\end{definition}

\begin{theorem}[插值误差]
设 $f \in C^{n+1}[a, b]$，$L_n$ 是 $f$ 在节点 $x_0, \ldots, x_n$ 处的 Lagrange 插值多项式，则：
$$f(x) - L_n(x) = \frac{f^{(n+1)}(\xi)}{(n+1)!} \prod_{i=0}^{n} (x - x_i)$$
其中 $\xi$ 在包含 $x, x_0, \ldots, x_n$ 的最小区间内。
\end{theorem}

\begin{definition}[Newton 插值]
使用差商定义的插值多项式：
$$N_n(x) = f[x_0] + f[x_0, x_1](x - x_0) + f[x_0, x_1, x_2](x - x_0)(x - x_1) + \cdots$$
其中差商递归定义为：
$$f[x_i] = f(x_i), \quad f[x_i, \ldots, x_{i+k}] = \frac{f[x_{i+1}, \ldots, x_{i+k}] - f[x_i, \ldots, x_{i+k-1}]}{x_{i+k} - x_i}$$
\end{definition}

\begin{definition}[样条插值]
三次样条 $S(x)$ 满足：
\begin{enumerate}
\item $S$ 在每个子区间 $[x_i, x_{i+1}]$ 上是三次多项式
\item $S(x_i) = y_i$
\item $S, S', S''$ 在 $(a, b)$ 内连续
\end{enumerate}

边界条件（常用）：
\begin{itemize}
\item 自然边界：$S''(a) = S''(b) = 0$
\item 固定边界：$S'(a) = f'(a)$，$S'(b) = f'(b)$
\end{itemize}
\end{definition}

\begin{theorem}[最佳平方逼近]
在内积空间中，$f$ 对子空间 $V$ 的最佳平方逼近是满足：
$$\|f - p^*\| = \min_{p \in V} \|f - p\|$$
的 $p^* \in V$。

若 $\{\phi_0, \ldots, \phi_n\}$ 是 $V$ 的基，则 $p^* = \sum_{i=0}^{n} c_i \phi_i$，其中系数满足正规方程：
$$Gc = b$$
$G_{ij} = \langle \phi_i, \phi_j \rangle$，$b_i = \langle f, \phi_i \rangle$。
\end{theorem}

\subsection{数值积分}

\begin{definition}[Newton-Cotes 公式]
基于等距节点的数值积分：
$$\int_a^b f(x) dx \approx \sum_{i=0}^{n} A_i f(x_i)$$
其中 $x_i = a + ih$，$h = (b-a)/n$。
\end{definition}

\begin{theorem}[常用求积公式]
\begin{enumerate}
\item 梯形公式：$\int_a^b f(x) dx \approx \frac{h}{2}[f(a) + f(b)]$，误差 $O(h^3f'')$
\item Simpson 公式：$\int_a^b f(x) dx \approx \frac{h}{6}[f(a) + 4f(\frac{a+b}{2}) + f(b)]$，误差 $O(h^5f^{(4)})$
\item 复合梯形公式：$\int_a^b f(x) dx \approx h[\frac{1}{2}f_0 + f_1 + \cdots + f_{n-1} + \frac{1}{2}f_n]$
\end{enumerate}
\end{theorem}

\begin{theorem}[Gauss 求积]
选取最优节点使得公式对尽可能高次多项式精确。$n$ 点 Gauss 求积对 $2n-1$ 次多项式精确。

Gauss-Legendre 公式（$[-1, 1]$ 上）：
$$\int_{-1}^{1} f(x) dx \approx \sum_{i=1}^{n} w_i f(x_i)$$
其中 $x_i$ 是 Legendre 多项式 $P_n(x)$ 的根。
\end{theorem}

\begin{theorem}[Romberg 积分]
使用 Richardson 外推提高精度。设 $T(h)$ 是步长 $h$ 的梯形公式结果，则：
$$T_1(h) = \frac{4T(h/2) - T(h)}{3}$$
是 Simpson 公式。继续外推得到更高精度。
\end{theorem}

\subsection{线性方程组数值解}

\begin{theorem}[Gauss 消元法]
将 $Ax = b$ 化为上三角形式 $Ux = c$，再回代求解。

时间复杂度：$O(n^3)$
\end{theorem}

\begin{theorem}[LU 分解]
若 $A$ 的所有顺序主子式非零，则存在唯一的 $LU$ 分解 $A = LU$，其中 $L$ 是单位下三角阵，$U$ 是上三角阵。

解 $Ax = b$：先解 $Ly = b$，再解 $Ux = y$。
\end{theorem}

\begin{theorem}[Cholesky 分解]
若 $A$ 是对称正定矩阵，则存在唯一的下三角阵 $L$（对角元为正）使得 $A = LL^T$。

计算量约为 $LU$ 分解的一半。
\end{theorem}

\begin{theorem}[迭代法]
将 $Ax = b$ 改写为 $x = Bx + c$ 的形式，迭代：
$$x^{(k+1)} = Bx^{(k)} + c$$

收敛条件：迭代矩阵 $B$ 的谱半径 $\rho(B) < 1$。
\end{theorem}

\begin{theorem}[Jacobi 迭代]
$$x_i^{(k+1)} = \frac{1}{a_{ii}}\left(b_i - \sum_{j \neq i} a_{ij}x_j^{(k)}\right)$$
\end{theorem}

\begin{theorem}[Gauss-Seidel 迭代]
$$x_i^{(k+1)} = \frac{1}{a_{ii}}\left(b_i - \sum_{j < i} a_{ij}x_j^{(k+1)} - \sum_{j > i} a_{ij}x_j^{(k)}\right)$$
\end{theorem}

\begin{theorem}[SOR 迭代]
$$x_i^{(k+1)} = (1-\omega)x_i^{(k)} + \frac{\omega}{a_{ii}}\left(b_i - \sum_{j < i} a_{ij}x_j^{(k+1)} - \sum_{j > i} a_{ij}x_j^{(k)}\right)$$
其中 $0 < \omega < 2$。
\end{theorem}

\begin{theorem}[共轭梯度法]
对于对称正定矩阵 $A$，共轭梯度法在精确算术下至多 $n$ 步收敛。

基本思想：在 Krylov 子空间 $\mathcal{K}_k(A, r_0) = \text{span}\{r_0, Ar_0, \ldots, A^{k-1}r_0\}$ 中寻找最优解。
\end{theorem}

% ================================
% 第八章：图论与组合
% ================================

\section{图论与组合数学}

\subsection{图的基本概念}

\begin{definition}[图]
图 $G = (V, E)$ 由顶点集 $V$ 和边集 $E$ 组成。边可以是有向的（有序对）或无向的（无序对）。

简单图：无自环、无重边。
多重图：允许重边。
\end{definition}

\begin{definition}[图的度]
顶点 $v$ 的度 $\deg(v)$ 是与 $v$ 关联的边数。

有向图中：入度 $\deg^-(v)$，出度 $\deg^+(v)$。
\end{definition}

\begin{theorem}[握手定理]
$$\sum_{v \in V} \deg(v) = 2|E|$$
推论：度为奇数的顶点个数是偶数。
\end{theorem}

\begin{definition}[路径与连通]
\begin{enumerate}
\item 路径：顶点序列 $v_0, v_1, \ldots, v_k$，相邻顶点之间有边
\item 简单路径：顶点不重复
\item 回路/圈：起点等于终点的路径
\item 连通：任意两点之间有路径
\item 连通分量：极大连通子图
\end{enumerate}
\end{definition}

\begin{definition}[特殊图]
\begin{enumerate}
\item 完全图 $K_n$：任意两点之间有边
\item 完全二部图 $K_{m,n}$：顶点分为两组，不同组之间有边
\item 圈 $C_n$：$n$ 个顶点的圈
\item 路径 $P_n$：$n$ 个顶点的路径
\item 星图 $S_n$：中心顶点连接 $n$ 个叶子
\end{enumerate}
\end{definition}

\subsection{树}

\begin{definition}[树]
树是无圈的连通图。

森林是无圈的图（可能不连通）。
\end{definition}

\begin{theorem}[树的等价定义]
设 $G$ 是 $n$ 个顶点的图。以下条件等价：
\begin{enumerate}
\item $G$ 是树
\item $G$ 连通且有 $n-1$ 条边
\item $G$ 无圈且有 $n-1$ 条边
\item 任意两点之间恰有一条路径
\item $G$ 连通，删除任一边则不连通
\item $G$ 无圈，添加任一边则有圈
\end{enumerate}
\end{theorem}

\begin{theorem}[Cayley 公式]
$n$ 个标记顶点的不同生成树数目是 $n^{n-2}$。
\end{theorem}

\begin{definition}[生成树]
图 $G$ 的生成树是包含所有顶点的树子图。
\end{definition}

\begin{theorem}[最小生成树]
给定带权连通图 $G$，最小生成树是权和最小的生成树。

算法：
\begin{enumerate}
\item Kruskal：按权从小到大选边，不形成圈则加入
\item Prim：从一个顶点开始，每次加入最小权边扩展
\end{enumerate}
\end{theorem}

\subsection{欧拉路与哈密顿路}

\begin{definition}[欧拉路与欧拉回路]
欧拉路：经过每条边恰好一次的路径。

欧拉回路：起点等于终点的欧拉路。
\end{definition}

\begin{theorem}[欧拉定理]
连通图 $G$ 有欧拉回路当且仅当每个顶点度为偶数。

连通图 $G$ 有欧拉路当且仅当恰有 0 或 2 个奇度顶点。
\end{theorem}

\begin{definition}[哈密顿路与哈密顿回路]
哈密顿路：经过每个顶点恰好一次的路径。

哈密顿回路：起点等于终点的哈密顿路。
\end{definition}

\begin{theorem}[Ore 定理]
设 $G$ 是 $n \geq 3$ 个顶点的简单图。若对于每对不相邻顶点 $u, v$，有 $\deg(u) + \deg(v) \geq n$，则 $G$ 有哈密顿回路。
\end{theorem}

\begin{theorem}[Dirac 定理]
设 $G$ 是 $n \geq 3$ 个顶点的简单图。若每个顶点度至少为 $n/2$，则 $G$ 有哈密顿回路。
\end{theorem}

\subsection{图着色}

\begin{definition}[顶点着色]
图 $G$ 的 $k$-着色是函数 $c: V \to \{1, 2, \ldots, k\}$，使得相邻顶点颜色不同。

色数 $\chi(G)$ 是使 $G$ 可着色的最小 $k$。
\end{definition}

\begin{theorem}[四色定理]
每个平面图的色数至多为 4。
\end{theorem}

\begin{theorem}[Brooks 定理]
设 $G$ 是连通简单图，不是完全图也不是奇圈。则 $\chi(G) \leq \Delta(G)$，其中 $\Delta(G)$ 是最大度。
\end{theorem}

\begin{definition}[边着色]
边着色使得共享顶点的边颜色不同。

边色数 $\chi'(G)$ 是最小边着色数。
\end{definition}

\begin{theorem}[Vizing 定理]
对于简单图 $G$：$\Delta(G) \leq \chi'(G) \leq \Delta(G) + 1$。
\end{theorem}

\subsection{组合计数}

\begin{theorem}[加法原理与乘法原理]
加法原理：若事件 $A$ 有 $m$ 种方式，事件 $B$ 有 $n$ 种方式，$A, B$ 互斥，则 $A$ 或 $B$ 有 $m + n$ 种方式。

乘法原理：若事件 $A$ 有 $m$ 种方式，对于 $A$ 的每种方式，事件 $B$ 有 $n$ 种方式，则 $A$ 后 $B$ 有 $mn$ 种方式。
\end{theorem}

\begin{theorem}[排列与组合]
从 $n$ 个不同元素中取 $k$ 个：

排列（有序）：$P(n, k) = \frac{n!}{(n-k)!}$

组合（无序）：$C(n, k) = \binom{n}{k} = \frac{n!}{k!(n-k)!}$
\end{theorem}

\begin{theorem}[二项式定理]
$$(x + y)^n = \sum_{k=0}^{n} \binom{n}{k} x^{n-k} y^k$$
\end{theorem}

\begin{theorem}[多项式定理]
$$(x_1 + x_2 + \cdots + x_m)^n = \sum_{k_1 + k_2 + \cdots + k_m = n} \binom{n}{k_1, k_2, \ldots, k_m} x_1^{k_1} x_2^{k_2} \cdots x_m^{k_m}$$
其中多项式系数 $\binom{n}{k_1, k_2, \ldots, k_m} = \frac{n!}{k_1! k_2! \cdots k_m!}$。
\end{theorem}

\begin{theorem}[容斥原理]
$$\left|\bigcup_{i=1}^{n} A_i\right| = \sum_{i} |A_i| - \sum_{i < j} |A_i \cap A_j| + \sum_{i < j < k} |A_i \cap A_j \cap A_k| - \cdots$$
\end{theorem}

\begin{theorem}[鸽巢原理]
若 $n+1$ 只鸽子放入 $n$ 个巢，则至少有一个巢有 $\geq 2$ 只鸽子。

一般形式：若 $n$ 个物体放入 $k$ 个盒子，则至少有一个盒子有 $\geq \lceil n/k \rceil$ 个物体。
\end{theorem}

\begin{definition}[生成函数]
序列 $\{a_n\}$ 的普通生成函数：$G(x) = \sum_{n=0}^{\infty} a_n x^n$

指数生成函数：$E(x) = \sum_{n=0}^{\infty} a_n \frac{x^n}{n!}$
\end{definition}

\begin{theorem}[Catalan 数]
Catalan 数 $C_n = \frac{1}{n+1}\binom{2n}{n}$ 计数许多组合对象，包括：
\begin{enumerate}
\item $n$ 对括号的合法匹配数
\item $n+1$ 个节点的不同二叉树数
\item 从 $(0,0)$ 到 $(n,n)$ 不越过对角线的格路数
\item $n+2$ 边形的三角剖分数
\end{enumerate}

递推关系：$C_0 = 1$，$C_{n+1} = \sum_{i=0}^{n} C_i C_{n-i}$

生成函数：$C(x) = \frac{1 - \sqrt{1 - 4x}}{2x}$
\end{theorem}

% ================================
% 结论
% ================================

\section{结论}

本文档涵盖了数学分析、线性代数、概率论与统计、复分析、拓扑学、常微分方程、数值分析以及图论与组合数学等多个数学领域的核心内容。这些内容构成了现代数学和计算机科学的理论基础。

通过本文档的测试，可以验证 Code Reader 在处理大型 LaTeX 文档时的性能表现，包括：

\begin{enumerate}
\item 增量解析效率：IncrementalLatexEngine 的增量解析性能
\item Tree-sitter 切换：大文件自动启用 Tree-sitter 模式的效果
\item 渲染性能：数学公式的 Web Worker 后台渲染效率
\item 滚动流畅度：IntersectionObserver 可视区域优先渲染的效果
\item 内存占用：LRU 缓存策略对内存使用的优化
\end{enumerate}

这些测试对于优化编辑器的用户体验至关重要。

% ================================
% 附录：更多数学公式
% ================================

\appendix

\section{常用数学公式汇编}

\subsection{三角函数公式}

\begin{align}
\sin^2\theta + \cos^2\theta &= 1 \\
1 + \tan^2\theta &= \sec^2\theta \\
1 + \cot^2\theta &= \csc^2\theta \\
\sin(\alpha \pm \beta) &= \sin\alpha\cos\beta \pm \cos\alpha\sin\beta \\
\cos(\alpha \pm \beta) &= \cos\alpha\cos\beta \mp \sin\alpha\sin\beta \\
\tan(\alpha \pm \beta) &= \frac{\tan\alpha \pm \tan\beta}{1 \mp \tan\alpha\tan\beta} \\
\sin 2\theta &= 2\sin\theta\cos\theta \\
\cos 2\theta &= \cos^2\theta - \sin^2\theta = 2\cos^2\theta - 1 = 1 - 2\sin^2\theta \\
\sin\alpha + \sin\beta &= 2\sin\frac{\alpha+\beta}{2}\cos\frac{\alpha-\beta}{2} \\
\cos\alpha + \cos\beta &= 2\cos\frac{\alpha+\beta}{2}\cos\frac{\alpha-\beta}{2} \\
\sin\alpha\cos\beta &= \frac{1}{2}[\sin(\alpha+\beta) + \sin(\alpha-\beta)] \\
\cos\alpha\cos\beta &= \frac{1}{2}[\cos(\alpha+\beta) + \cos(\alpha-\beta)]
\end{align}

\subsection{常用积分表}

\begin{align}
\int \frac{1}{x^2 + a^2} dx &= \frac{1}{a}\arctan\frac{x}{a} + C \\
\int \frac{1}{x^2 - a^2} dx &= \frac{1}{2a}\ln\left|\frac{x-a}{x+a}\right| + C \\
\int \frac{1}{\sqrt{a^2 - x^2}} dx &= \arcsin\frac{x}{a} + C \\
\int \frac{1}{\sqrt{x^2 + a^2}} dx &= \ln(x + \sqrt{x^2 + a^2}) + C \\
\int \frac{1}{\sqrt{x^2 - a^2}} dx &= \ln|x + \sqrt{x^2 - a^2}| + C \\
\int \sqrt{a^2 - x^2} dx &= \frac{x}{2}\sqrt{a^2 - x^2} + \frac{a^2}{2}\arcsin\frac{x}{a} + C \\
\int \sqrt{x^2 \pm a^2} dx &= \frac{x}{2}\sqrt{x^2 \pm a^2} \pm \frac{a^2}{2}\ln|x + \sqrt{x^2 \pm a^2}| + C \\
\int \tan x dx &= -\ln|\cos x| + C \\
\int \cot x dx &= \ln|\sin x| + C \\
\int \sec x dx &= \ln|\sec x + \tan x| + C \\
\int \csc x dx &= \ln|\csc x - \cot x| + C \\
\int e^{ax}\sin bx dx &= \frac{e^{ax}}{a^2 + b^2}(a\sin bx - b\cos bx) + C \\
\int e^{ax}\cos bx dx &= \frac{e^{ax}}{a^2 + b^2}(a\cos bx + b\sin bx) + C \\
\int x^n e^x dx &= x^n e^x - n\int x^{n-1} e^x dx \\
\int \ln x dx &= x\ln x - x + C \\
\int x^n \ln x dx &= \frac{x^{n+1}}{n+1}\ln x - \frac{x^{n+1}}{(n+1)^2} + C \quad (n \neq -1)
\end{align}

\subsection{常用级数}

\begin{align}
\sum_{k=1}^{n} k &= \frac{n(n+1)}{2} \\
\sum_{k=1}^{n} k^2 &= \frac{n(n+1)(2n+1)}{6} \\
\sum_{k=1}^{n} k^3 &= \left[\frac{n(n+1)}{2}\right]^2 \\
\sum_{k=0}^{n} x^k &= \frac{1 - x^{n+1}}{1 - x} \quad (x \neq 1) \\
\sum_{k=0}^{\infty} x^k &= \frac{1}{1 - x} \quad (|x| < 1) \\
\sum_{k=1}^{\infty} \frac{1}{k^2} &= \frac{\pi^2}{6} \\
\sum_{k=1}^{\infty} \frac{1}{k^4} &= \frac{\pi^4}{90} \\
\sum_{k=1}^{\infty} \frac{(-1)^{k+1}}{k} &= \ln 2 \\
\sum_{k=0}^{\infty} \frac{(-1)^k}{2k+1} &= \frac{\pi}{4} \\
\sum_{k=1}^{\infty} \frac{1}{k(k+1)} &= 1 \\
\sum_{k=1}^{\infty} \frac{1}{k!} &= e - 1
\end{align}

\subsection{Gamma 函数与 Beta 函数}

\begin{definition}[Gamma 函数]
$$\Gamma(z) = \int_0^{\infty} t^{z-1} e^{-t} dt, \quad \text{Re}(z) > 0$$

性质：
\begin{align}
\Gamma(n+1) &= n! \quad (n \in \mathbb{N}) \\
\Gamma(z+1) &= z\Gamma(z) \\
\Gamma(1/2) &= \sqrt{\pi} \\
\Gamma(z)\Gamma(1-z) &= \frac{\pi}{\sin(\pi z)}
\end{align}
\end{definition}

\begin{definition}[Beta 函数]
$$B(x, y) = \int_0^1 t^{x-1}(1-t)^{y-1} dt = \frac{\Gamma(x)\Gamma(y)}{\Gamma(x+y)}$$
\end{definition}

\subsection{Fourier 级数}

\begin{theorem}[Fourier 级数展开]
设 $f$ 在 $[-\pi, \pi]$ 上可积，其 Fourier 级数为：
$$f(x) \sim \frac{a_0}{2} + \sum_{n=1}^{\infty} (a_n \cos nx + b_n \sin nx)$$
其中：
\begin{align}
a_n &= \frac{1}{\pi} \int_{-\pi}^{\pi} f(x) \cos nx \, dx \\
b_n &= \frac{1}{\pi} \int_{-\pi}^{\pi} f(x) \sin nx \, dx
\end{align}
\end{theorem}

\begin{theorem}[Parseval 等式]
$$\frac{1}{\pi} \int_{-\pi}^{\pi} |f(x)|^2 dx = \frac{|a_0|^2}{2} + \sum_{n=1}^{\infty} (|a_n|^2 + |b_n|^2)$$
\end{theorem}

常用 Fourier 级数展开：
\begin{align}
x &= 2\sum_{n=1}^{\infty} \frac{(-1)^{n+1}}{n} \sin nx \quad (-\pi < x < \pi) \\
|x| &= \frac{\pi}{2} - \frac{4}{\pi}\sum_{k=0}^{\infty} \frac{\cos(2k+1)x}{(2k+1)^2} \quad (-\pi \leq x \leq \pi) \\
x^2 &= \frac{\pi^2}{3} + 4\sum_{n=1}^{\infty} \frac{(-1)^n}{n^2} \cos nx \quad (-\pi \leq x \leq \pi)
\end{align}

\section{更多定理与公式}

\subsection{测度论基础}

\begin{definition}[$\sigma$-代数]
设 $X$ 是集合，$\mathcal{F} \subseteq \mathcal{P}(X)$ 称为 $X$ 上的 $\sigma$-代数，若：
\begin{enumerate}
\item $X \in \mathcal{F}$
\item $A \in \mathcal{F} \Rightarrow A^c \in \mathcal{F}$
\item $\{A_n\}_{n=1}^{\infty} \subseteq \mathcal{F} \Rightarrow \bigcup_{n=1}^{\infty} A_n \in \mathcal{F}$
\end{enumerate}
\end{definition}

\begin{definition}[测度]
设 $(X, \mathcal{F})$ 是可测空间。函数 $\mu: \mathcal{F} \to [0, \infty]$ 称为测度，若：
\begin{enumerate}
\item $\mu(\emptyset) = 0$
\item 可数可加性：对于两两不相交的 $\{A_n\} \subseteq \mathcal{F}$，
$$\mu\left(\bigcup_{n=1}^{\infty} A_n\right) = \sum_{n=1}^{\infty} \mu(A_n)$$
\end{enumerate}
\end{definition}

\begin{theorem}[Lebesgue 测度的性质]
设 $\lambda$ 是 $\mathbb{R}^n$ 上的 Lebesgue 测度：
\begin{enumerate}
\item 区间 $[a, b]$ 的测度为 $b - a$
\item $\lambda$ 是平移不变的：$\lambda(A + x) = \lambda(A)$
\item $\lambda$ 是完备的
\item 每个可数集的测度为 0
\item 存在不可测集（Vitali 集）
\end{enumerate}
\end{theorem}

\begin{theorem}[单调收敛定理]
设 $f_n \geq 0$ 可测，$f_n \nearrow f$（单调递增收敛到 $f$），则：
$$\lim_{n \to \infty} \int f_n d\mu = \int f d\mu$$
\end{theorem}

\begin{theorem}[Fatou 引理]
设 $f_n \geq 0$ 可测，则：
$$\int \liminf_{n \to \infty} f_n d\mu \leq \liminf_{n \to \infty} \int f_n d\mu$$
\end{theorem}

\begin{theorem}[控制收敛定理]
设 $f_n \to f$ 逐点收敛，存在可积函数 $g$ 使得 $|f_n| \leq g$，则：
$$\lim_{n \to \infty} \int f_n d\mu = \int f d\mu$$
\end{theorem}

\subsection{泛函分析}

\begin{definition}[范数空间]
设 $V$ 是向量空间。函数 $\|\cdot\|: V \to [0, \infty)$ 称为范数，若：
\begin{enumerate}
\item $\|x\| = 0 \Leftrightarrow x = 0$
\item $\|\alpha x\| = |\alpha| \|x\|$
\item $\|x + y\| \leq \|x\| + \|y\|$（三角不等式）
\end{enumerate}
\end{definition}

\begin{definition}[Banach 空间]
完备的范数空间称为 Banach 空间。
\end{definition}

\begin{definition}[内积空间]
设 $V$ 是复向量空间。函数 $\langle\cdot, \cdot\rangle: V \times V \to \mathbb{C}$ 称为内积，若：
\begin{enumerate}
\item $\langle x, x \rangle \geq 0$，等号当且仅当 $x = 0$
\item $\langle x, y \rangle = \overline{\langle y, x \rangle}$
\item $\langle \alpha x + \beta y, z \rangle = \alpha\langle x, z \rangle + \beta\langle y, z \rangle$
\end{enumerate}
\end{definition}

\begin{definition}[Hilbert 空间]
完备的内积空间称为 Hilbert 空间。
\end{definition}

\begin{theorem}[Riesz 表示定理]
设 $H$ 是 Hilbert 空间，$\varphi: H \to \mathbb{C}$ 是有界线性泛函。则存在唯一的 $y \in H$ 使得：
$$\varphi(x) = \langle x, y \rangle, \quad \forall x \in H$$
且 $\|\varphi\| = \|y\|$。
\end{theorem}

\begin{theorem}[正交投影定理]
设 $H$ 是 Hilbert 空间，$M \subseteq H$ 是闭子空间。则对于每个 $x \in H$，存在唯一的 $m \in M$ 使得：
$$\|x - m\| = \inf_{y \in M} \|x - y\|$$
且 $x - m \perp M$。
\end{theorem}

\begin{theorem}[Hahn-Banach 定理]
设 $X$ 是范数空间，$M \subseteq X$ 是子空间，$f: M \to \mathbb{C}$ 是有界线性泛函。则存在有界线性泛函 $F: X \to \mathbb{C}$ 使得：
\begin{enumerate}
\item $F|_M = f$
\item $\|F\| = \|f\|$
\end{enumerate}
\end{theorem}

\begin{theorem}[开映射定理]
设 $X, Y$ 是 Banach 空间，$T: X \to Y$ 是有界线性算子且满射。则 $T$ 是开映射（将开集映为开集）。
\end{theorem}

\begin{theorem}[闭图像定理]
设 $X, Y$ 是 Banach 空间，$T: X \to Y$ 是线性算子。$T$ 有界当且仅当 $T$ 的图像 $\{(x, Tx) : x \in X\}$ 在 $X \times Y$ 中是闭的。
\end{theorem}

\begin{theorem}[一致有界性原理（Banach-Steinhaus 定理）]
设 $X$ 是 Banach 空间，$Y$ 是范数空间，$\{T_\alpha\}$ 是 $X$ 到 $Y$ 的有界线性算子族。若对于每个 $x \in X$，$\sup_\alpha \|T_\alpha x\| < \infty$，则 $\sup_\alpha \|T_\alpha\| < \infty$。
\end{theorem}

\subsection{偏微分方程}

\begin{definition}[偏微分方程分类]
二阶线性 PDE：
$$a u_{xx} + 2b u_{xy} + c u_{yy} + \text{低阶项} = 0$$

判别式 $\Delta = b^2 - ac$：
\begin{enumerate}
\item $\Delta > 0$：双曲型（如波动方程）
\item $\Delta = 0$：抛物型（如热传导方程）
\item $\Delta < 0$：椭圆型（如 Laplace 方程）
\end{enumerate}
\end{definition}

\begin{theorem}[热传导方程]
$$u_t = k u_{xx}$$

初边值问题的解（$0 < x < L$）：
$$u(x, t) = \sum_{n=1}^{\infty} B_n \sin\frac{n\pi x}{L} e^{-k(n\pi/L)^2 t}$$
其中 $B_n$ 由初始条件确定。

无界域上的解（基本解）：
$$u(x, t) = \frac{1}{\sqrt{4\pi kt}} \int_{-\infty}^{\infty} f(\xi) e^{-(x-\xi)^2/(4kt)} d\xi$$
\end{theorem}

\begin{theorem}[波动方程]
$$u_{tt} = c^2 u_{xx}$$

d'Alembert 解（无界域）：
$$u(x, t) = \frac{1}{2}[f(x - ct) + f(x + ct)] + \frac{1}{2c}\int_{x-ct}^{x+ct} g(\xi) d\xi$$
其中 $u(x, 0) = f(x)$，$u_t(x, 0) = g(x)$。

有界域解（$0 < x < L$）：
$$u(x, t) = \sum_{n=1}^{\infty} \left(A_n \cos\frac{n\pi ct}{L} + B_n \sin\frac{n\pi ct}{L}\right) \sin\frac{n\pi x}{L}$$
\end{theorem}

\begin{theorem}[Laplace 方程]
$$\nabla^2 u = u_{xx} + u_{yy} = 0$$

满足 Laplace 方程的函数称为调和函数。

\textbf{平均值性质：}调和函数在任一点的值等于该点任意球面上的平均值。

\textbf{最大值原理：}非常数调和函数在区域内部不取得最大值或最小值。

\textbf{Poisson 公式（圆域）：}
$$u(r, \theta) = \frac{1}{2\pi} \int_0^{2\pi} \frac{R^2 - r^2}{R^2 - 2Rr\cos(\theta - \phi) + r^2} f(\phi) d\phi$$
\end{theorem}

\begin{theorem}[Green 函数]
Laplace 方程的 Green 函数 $G(x, \xi)$ 满足：
$$-\nabla^2 G = \delta(x - \xi)$$
边界条件：$G|_{\partial\Omega} = 0$。

解的表示：
$$u(x) = \int_{\partial\Omega} f(\xi) \frac{\partial G}{\partial n_\xi} dS_\xi$$
其中 $f$ 是边界条件。
\end{theorem}

\subsection{变分法}

\begin{definition}[泛函]
泛函是从函数空间到实数的映射，如：
$$J[y] = \int_a^b F(x, y, y') dx$$
\end{definition}

\begin{theorem}[欧拉-拉格朗日方程]
泛函 $J[y] = \int_a^b F(x, y, y') dx$ 取极值的必要条件是：
$$\frac{\partial F}{\partial y} - \frac{d}{dx}\frac{\partial F}{\partial y'} = 0$$
\end{theorem}

\begin{example}[最短路径问题]
两点间最短路径：$J[y] = \int_a^b \sqrt{1 + (y')^2} dx$

欧拉方程给出 $y'' = 0$，即直线。
\end{example}

\begin{example}[最速降线问题]
质点在重力场中从 $A$ 点滑到 $B$ 点的最短时间路径。

泛函：$T[y] = \int_0^{x_1} \sqrt{\frac{1 + (y')^2}{2gy}} dx$

解：摆线（cycloid）
$$x = a(\theta - \sin\theta), \quad y = a(1 - \cos\theta)$$
\end{example}

\begin{theorem}[等周问题]
在周长固定的条件下求面积最大的曲线。

使用 Lagrange 乘数法：
$$J[y] = \int_{-a}^{a} y dx - \lambda \int_{-a}^{a} \sqrt{1 + (y')^2} dx$$

解：圆
\end{theorem}

\subsection{复杂系统与混沌}

\begin{definition}[动力系统]
离散动力系统：$x_{n+1} = f(x_n)$

连续动力系统：$\dot{x} = f(x)$

其中 $x \in \mathbb{R}^n$，$f: \mathbb{R}^n \to \mathbb{R}^n$。
\end{definition}

\begin{definition}[不动点与周期轨道]
不动点：$f(x^*) = x^*$

周期 $p$ 轨道：$f^p(x) = x$ 但 $f^k(x) \neq x$（$k < p$）
\end{definition}

\begin{definition}[稳定性]
不动点 $x^*$ 的稳定性由 Jacobi 矩阵 $Df(x^*)$ 的特征值决定：
\begin{enumerate}
\item 渐近稳定：所有特征值模小于 1（离散）或实部小于 0（连续）
\item 不稳定：存在特征值模大于 1 或实部大于 0
\end{enumerate}
\end{definition}

\begin{definition}[混沌]
混沌系统的特征：
\begin{enumerate}
\item 对初值敏感（Lyapunov 指数为正）
\item 拓扑传递性
\item 稠密周期轨道
\end{enumerate}
\end{definition}

\begin{example}[Logistic 映射]
$$x_{n+1} = rx_n(1 - x_n), \quad r \in [0, 4]$$

\begin{itemize}
\item $r < 1$：收敛到 0
\item $1 < r < 3$：收敛到不动点 $(r-1)/r$
\item $3 < r < 3.57$：周期倍化分岔
\item $r > 3.57$：混沌（存在周期窗口）
\end{itemize}
\end{example}

\begin{example}[Lorenz 系统]
$$\begin{cases}
\dot{x} = \sigma(y - x) \\
\dot{y} = x(\rho - z) - y \\
\dot{z} = xy - \beta z
\end{cases}$$

经典参数：$\sigma = 10$，$\rho = 28$，$\beta = 8/3$。

产生著名的 Lorenz 吸引子（蝴蝶形状）。
\end{example}

\begin{theorem}[Lyapunov 指数]
Lyapunov 指数衡量相邻轨道的分离速率：
$$\lambda = \lim_{n \to \infty} \frac{1}{n} \sum_{k=0}^{n-1} \ln|f'(x_k)|$$

$\lambda > 0$ 表示混沌。
\end{theorem}

\subsection{信息论}

\begin{definition}[熵]
离散随机变量 $X$ 的熵：
$$H(X) = -\sum_{x} P(X = x) \log_2 P(X = x)$$

熵度量不确定性或平均信息量。
\end{definition}

\begin{theorem}[熵的性质]
\begin{enumerate}
\item $H(X) \geq 0$，等号当且仅当 $X$ 确定
\item $H(X) \leq \log_2 n$（$n$ 个可能值），等号当且仅当均匀分布
\item 条件熵：$H(X|Y) = \sum_y P(Y = y) H(X|Y = y)$
\item 链式法则：$H(X, Y) = H(X) + H(Y|X)$
\item $H(X|Y) \leq H(X)$，等号当且仅当 $X, Y$ 独立
\end{enumerate}
\end{theorem}

\begin{definition}[互信息]
$$I(X; Y) = H(X) - H(X|Y) = H(Y) - H(Y|X) = H(X) + H(Y) - H(X, Y)$$

互信息度量两个变量的相互依赖程度。
\end{definition}

\begin{definition}[相对熵（KL 散度）]
$$D_{KL}(P \| Q) = \sum_x P(x) \log \frac{P(x)}{Q(x)}$$

$D_{KL}(P \| Q) \geq 0$，等号当且仅当 $P = Q$。
\end{definition}

\begin{theorem}[信道容量]
信道容量是互信息的最大值：
$$C = \max_{P(X)} I(X; Y)$$

Shannon 信道编码定理：信息传输速率可以任意接近信道容量而误码率任意小。
\end{theorem}

\begin{theorem}[数据压缩基本定理]
i.i.d. 信源的最优无损压缩率为熵 $H(X)$。
\end{theorem}

\subsection{优化理论}

\begin{definition}[凸优化]
凸优化问题：
$$\min_{x} f(x) \quad \text{s.t.} \quad g_i(x) \leq 0, \quad h_j(x) = 0$$
其中 $f, g_i$ 是凸函数，$h_j$ 是仿射函数。
\end{definition}

\begin{theorem}[KKT 条件]
设 $x^*$ 是凸优化问题的最优解，则存在 $\lambda_i^* \geq 0$，$\nu_j^*$ 使得：
\begin{align}
\nabla f(x^*) + \sum_i \lambda_i^* \nabla g_i(x^*) + \sum_j \nu_j^* \nabla h_j(x^*) &= 0 \\
\lambda_i^* g_i(x^*) &= 0 \quad \forall i \\
g_i(x^*) &\leq 0 \quad \forall i \\
h_j(x^*) &= 0 \quad \forall j
\end{align}
\end{theorem}

\begin{theorem}[强对偶性]
对于凸优化问题，若 Slater 条件成立（存在严格可行点），则强对偶性成立：
$$\min_x \max_{\lambda \geq 0, \nu} L(x, \lambda, \nu) = \max_{\lambda \geq 0, \nu} \min_x L(x, \lambda, \nu)$$
其中 Lagrangian：$L(x, \lambda, \nu) = f(x) + \sum_i \lambda_i g_i(x) + \sum_j \nu_j h_j(x)$。
\end{theorem}

\begin{theorem}[梯度下降法]
迭代：$x^{(k+1)} = x^{(k)} - \alpha_k \nabla f(x^{(k)})$

对于 $L$-Lipschitz 连续梯度的凸函数，步长 $\alpha = 1/L$ 保证收敛：
$$f(x^{(k)}) - f(x^*) \leq \frac{L\|x^{(0)} - x^*\|^2}{2k}$$
\end{theorem}

\begin{theorem}[牛顿法（优化）]
迭代：$x^{(k+1)} = x^{(k)} - [\nabla^2 f(x^{(k)})]^{-1} \nabla f(x^{(k)})$

二次收敛：$\|x^{(k+1)} - x^*\| \leq C\|x^{(k)} - x^*\|^2$
\end{theorem}

\begin{theorem}[线性规划]
$$\min c^T x \quad \text{s.t.} \quad Ax = b, \quad x \geq 0$$

单纯形法：在可行域的顶点之间移动，每步改进目标函数。

内点法：在可行域内部移动，多项式时间复杂度。
\end{theorem}

\subsection{机器学习数学基础}

\begin{definition}[监督学习]
给定训练数据 $\{(x_i, y_i)\}_{i=1}^n$，学习映射 $f: X \to Y$。

损失函数：$L(y, f(x))$，如：
\begin{itemize}
\item 平方损失：$(y - f(x))^2$
\item 交叉熵损失：$-y\log f(x) - (1-y)\log(1-f(x))$
\end{itemize}
\end{definition}

\begin{theorem}[偏差-方差分解]
对于均方误差：
$$\mathbb{E}[(y - \hat{f}(x))^2] = \text{Bias}^2[\hat{f}(x)] + \text{Var}[\hat{f}(x)] + \sigma^2$$
其中 $\sigma^2$ 是不可约误差。
\end{theorem}

\begin{definition}[正则化]
添加惩罚项防止过拟合：
$$\min_w \sum_i L(y_i, f_w(x_i)) + \lambda R(w)$$

常用正则化：
\begin{itemize}
\item L2（Ridge）：$R(w) = \|w\|_2^2$
\item L1（Lasso）：$R(w) = \|w\|_1$
\item 弹性网络：$R(w) = \alpha\|w\|_1 + (1-\alpha)\|w\|_2^2$
\end{itemize}
\end{definition}

\begin{theorem}[核方法]
核函数 $K(x, x') = \langle \phi(x), \phi(x') \rangle$ 允许在高维特征空间中隐式计算。

常用核：
\begin{itemize}
\item 多项式核：$K(x, x') = (x^T x' + c)^d$
\item RBF 核：$K(x, x') = \exp(-\gamma\|x - x'\|^2)$
\end{itemize}

Mercer 定理：$K$ 是有效核当且仅当 Gram 矩阵 $(K(x_i, x_j))$ 半正定。
\end{theorem}

\begin{theorem}[VC 维与泛化界]
假设类 $\mathcal{H}$ 的 VC 维 $d$ 衡量其复杂度。

泛化界：
$$\mathbb{E}[L(h)] \leq \hat{L}(h) + O\left(\sqrt{\frac{d \log n + \log(1/\delta)}{n}}\right)$$
以概率至少 $1 - \delta$ 成立，其中 $\hat{L}$ 是经验损失。
\end{theorem}

\begin{theorem}[反向传播]
对于神经网络 $f = f_L \circ f_{L-1} \circ \cdots \circ f_1$，链式法则给出：
$$\frac{\partial L}{\partial W_l} = \frac{\partial L}{\partial a_l} \cdot \frac{\partial a_l}{\partial W_l}$$
其中 $\frac{\partial L}{\partial a_l}$ 通过反向传播计算。
\end{theorem}

\section{数学符号速查}

\subsection{集合论符号}
\begin{align*}
\in, \notin &\quad \text{属于、不属于} \\
\subseteq, \subset &\quad \text{子集、真子集} \\
\cup, \cap &\quad \text{并集、交集} \\
\setminus, \triangle &\quad \text{差集、对称差} \\
\emptyset &\quad \text{空集} \\
\mathcal{P}(A) &\quad \text{幂集} \\
|A|, \#A &\quad \text{基数}
\end{align*}

\subsection{逻辑符号}
\begin{align*}
\land, \lor, \lnot &\quad \text{与、或、非} \\
\Rightarrow, \Leftrightarrow &\quad \text{蕴含、当且仅当} \\
\forall, \exists &\quad \text{对所有、存在} \\
\therefore, \because &\quad \text{所以、因为}
\end{align*}

\subsection{数集符号}
\begin{align*}
\mathbb{N}, \mathbb{Z}, \mathbb{Q} &\quad \text{自然数、整数、有理数} \\
\mathbb{R}, \mathbb{C} &\quad \text{实数、复数} \\
\mathbb{R}^n &\quad \text{$n$ 维实向量空间} \\
\mathbb{Z}_n, \mathbb{Z}/n\mathbb{Z} &\quad \text{模 $n$ 整数}
\end{align*}

\subsection{分析符号}
\begin{align*}
\lim, \limsup, \liminf &\quad \text{极限、上极限、下极限} \\
\sum, \prod &\quad \text{求和、求积} \\
\int, \oint &\quad \text{积分、环路积分} \\
\frac{d}{dx}, \frac{\partial}{\partial x} &\quad \text{导数、偏导数} \\
\nabla, \nabla^2 &\quad \text{梯度、拉普拉斯算子} \\
\|x\|, |x| &\quad \text{范数、绝对值/模}
\end{align*}

\subsection{代数符号}
\begin{align*}
\det, \text{tr} &\quad \text{行列式、迹} \\
\dim, \text{rank} &\quad \text{维数、秩} \\
\ker, \text{Im} &\quad \text{核、像} \\
\langle \cdot, \cdot \rangle &\quad \text{内积} \\
\oplus, \otimes &\quad \text{直和、张量积} \\
\cong, \simeq &\quad \text{同构、同伦}
\end{align*}

\subsection{概率符号}
\begin{align*}
P(A), \mathbb{P}(A) &\quad \text{概率} \\
E[X], \mathbb{E}[X] &\quad \text{期望} \\
\text{Var}(X), \sigma^2 &\quad \text{方差} \\
\text{Cov}(X, Y), \rho &\quad \text{协方差、相关系数} \\
X \sim D &\quad \text{$X$ 服从分布 $D$} \\
X \perp Y &\quad \text{$X, Y$ 独立}
\end{align*}

\end{document}

% ================================
% 扩展内容：更多数学定理
% ================================

\section{扩展内容}

\subsection{代数拓扑}

\begin{definition}[同伦]
设 $f, g: X \to Y$ 是连续映射。$f$ 和 $g$ 称为同伦的（记为 $f \simeq g$），若存在连续映射 $H: X \times [0, 1] \to Y$ 使得 $H(x, 0) = f(x)$，$H(x, 1) = g(x)$。
\end{definition}

\begin{definition}[基本群]
设 $X$ 是拓扑空间，$x_0 \in X$。$X$ 在 $x_0$ 处的基本群 $\pi_1(X, x_0)$ 由以 $x_0$ 为基点的回路的同伦等价类构成，运算为路径连接。
\end{definition}

\begin{theorem}[覆叠空间理论]
设 $p: \tilde{X} \to X$ 是覆叠映射。则：
\begin{enumerate}
\item 唯一路径提升：给定 $\tilde{X}$ 中的点和 $X$ 中的路径，存在唯一的路径提升
\item $\pi_1(\tilde{X}, \tilde{x}_0)$ 是 $\pi_1(X, x_0)$ 的子群
\item 如果 $\tilde{X}$ 是单连通的（万有覆叠空间），则 $\pi_1(X, x_0) \cong$ 覆叠变换群
\end{enumerate}
\end{theorem}

\begin{theorem}[Van Kampen 定理]
设 $X = U \cup V$，$U, V, U \cap V$ 都是道路连通的，$x_0 \in U \cap V$。则：
$$\pi_1(X, x_0) \cong \pi_1(U, x_0) *_{\pi_1(U \cap V, x_0)} \pi_1(V, x_0)$$
即基本群是 $\pi_1(U)$ 和 $\pi_1(V)$ 在 $\pi_1(U \cap V)$ 上的融合积。
\end{theorem}

\begin{definition}[同调群]
链复形 $\cdots \to C_{n+1} \xrightarrow{\partial_{n+1}} C_n \xrightarrow{\partial_n} C_{n-1} \to \cdots$ 的第 $n$ 个同调群：
$$H_n = \ker \partial_n / \text{Im } \partial_{n+1}$$
\end{definition}

\begin{theorem}[Eilenberg-Steenrod 公理]
同调论满足：
\begin{enumerate}
\item 同伦公理：同伦映射诱导相同的同调同态
\item 正合性公理：对 $(X, A)$ 有长正合序列
\item 切除公理：切除不影响相对同调
\item 维数公理：点的同调只在 $n = 0$ 时非零
\end{enumerate}
\end{theorem}

\subsection{黎曼几何}

\begin{definition}[黎曼度量]
流形 $M$ 上的黎曼度量是光滑的正定内积场 $g = (g_{ij})$，使得每个切空间 $T_p M$ 成为内积空间。

线元素：$ds^2 = \sum_{i,j} g_{ij} dx^i dx^j$
\end{definition}

\begin{definition}[联络与协变导数]
Levi-Civita 联络的 Christoffel 符号：
$$\Gamma^k_{ij} = \frac{1}{2} g^{kl} \left( \frac{\partial g_{il}}{\partial x^j} + \frac{\partial g_{jl}}{\partial x^i} - \frac{\partial g_{ij}}{\partial x^l} \right)$$

协变导数：
$$\nabla_{\partial_i} \partial_j = \Gamma^k_{ij} \partial_k$$
\end{definition}

\begin{definition}[曲率张量]
黎曼曲率张量：
$$R^l_{ijk} = \frac{\partial \Gamma^l_{jk}}{\partial x^i} - \frac{\partial \Gamma^l_{ik}}{\partial x^j} + \Gamma^m_{jk}\Gamma^l_{im} - \Gamma^m_{ik}\Gamma^l_{jm}$$

Ricci 张量：$R_{ij} = R^k_{ikj}$

标量曲率：$R = g^{ij} R_{ij}$
\end{definition}

\begin{theorem}[Gauss-Bonnet 定理]
设 $M$ 是紧致的 2 维黎曼流形，则：
$$\int_M K \, dA = 2\pi \chi(M)$$
其中 $K$ 是高斯曲率，$\chi(M)$ 是欧拉示性数。
\end{theorem}

\subsection{李群与李代数}

\begin{definition}[李群]
李群是同时具有光滑流形结构和群结构的数学对象，且群运算（乘法和逆）是光滑的。
\end{definition}

\begin{example}
常见李群：
\begin{itemize}
\item $GL(n, \mathbb{R})$：$n$ 阶可逆实矩阵
\item $SL(n, \mathbb{R})$：行列式为 1 的矩阵
\item $O(n)$：正交矩阵
\item $SO(n)$：行列式为 1 的正交矩阵
\item $U(n)$：酉矩阵
\item $SU(n)$：行列式为 1 的酉矩阵
\end{itemize}
\end{example}

\begin{definition}[李代数]
李群 $G$ 的李代数 $\mathfrak{g}$ 是单位元处的切空间 $T_e G$，配备李括号 $[\cdot, \cdot]$。

矩阵李群的李括号：$[A, B] = AB - BA$
\end{definition}

\begin{theorem}[指数映射]
$\exp: \mathfrak{g} \to G$ 由以下定义：
$$\exp(X) = e^X = \sum_{n=0}^{\infty} \frac{X^n}{n!}$$

对于矩阵李群，这就是矩阵指数。

$\exp$ 在 0 附近是局部微分同胚。
\end{theorem}

\begin{theorem}[Baker-Campbell-Hausdorff 公式]
$$\log(\exp X \exp Y) = X + Y + \frac{1}{2}[X, Y] + \frac{1}{12}[X, [X, Y]] - \frac{1}{12}[Y, [X, Y]] + \cdots$$
\end{theorem}

\subsection{表示论}

\begin{definition}[表示]
群 $G$ 的表示是同态 $\rho: G \to GL(V)$，其中 $V$ 是向量空间。

表示的维数是 $\dim V$。
\end{definition}

\begin{definition}[不可约表示]
表示 $(V, \rho)$ 称为不可约的，若 $V$ 没有非平凡的 $G$-不变子空间。
\end{definition}

\begin{theorem}[Schur 引理]
设 $\rho_1, \rho_2$ 是群 $G$ 的不可约表示。
\begin{enumerate}
\item 若 $\rho_1, \rho_2$ 不等价，则它们之间的任何 $G$-态射都是零
\item 若 $\rho$ 是复不可约表示的 $G$-自态射，则 $\phi = \lambda \cdot \text{id}$
\end{enumerate}
\end{theorem}

\begin{theorem}[Maschke 定理]
设 $G$ 是有限群，$|G|$ 与域的特征互素。则 $G$ 的每个表示都完全可约（可分解为不可约表示的直和）。
\end{theorem}

\begin{theorem}[特征理论]
表示的特征 $\chi(g) = \text{tr}(\rho(g))$ 是类函数（在共轭类上常值）。

正交关系：
$$\frac{1}{|G|} \sum_{g \in G} \chi_i(g) \overline{\chi_j(g)} = \delta_{ij}$$
\end{theorem}

\subsection{范畴论}

\begin{definition}[范畴]
范畴 $\mathcal{C}$ 由以下数据组成：
\begin{itemize}
\item 对象类 $\text{Ob}(\mathcal{C})$
\item 对于每对对象 $A, B$，态射集 $\text{Hom}(A, B)$
\item 态射的复合运算
\item 每个对象的恒等态射
\end{itemize}
满足结合律和单位元律。
\end{definition}

\begin{example}
常见范畴：
\begin{itemize}
\item $\mathbf{Set}$：集合与函数
\item $\mathbf{Grp}$：群与群同态
\item $\mathbf{Top}$：拓扑空间与连续映射
\item $\mathbf{Vect}_k$：$k$-向量空间与线性映射
\end{itemize}
\end{example}

\begin{definition}[函子]
函子 $F: \mathcal{C} \to \mathcal{D}$ 将对象和态射从一个范畴映射到另一个，保持复合和恒等。

共变函子：$F(g \circ f) = F(g) \circ F(f)$
反变函子：$F(g \circ f) = F(f) \circ F(g)$
\end{definition}

\begin{definition}[自然变换]
函子 $F, G: \mathcal{C} \to \mathcal{D}$ 之间的自然变换 $\eta$ 是态射族 $\{\eta_A: F(A) \to G(A)\}_{A \in \mathcal{C}}$，使得对于每个态射 $f: A \to B$，以下图表交换：
$$G(f) \circ \eta_A = \eta_B \circ F(f)$$
\end{definition}

\begin{definition}[伴随函子]
函子 $F: \mathcal{C} \to \mathcal{D}$ 和 $G: \mathcal{D} \to \mathcal{C}$ 称为伴随的（$F \dashv G$），若存在自然双射：
$$\text{Hom}_\mathcal{D}(F(A), B) \cong \text{Hom}_\mathcal{C}(A, G(B))$$
\end{definition}

\begin{theorem}[Yoneda 引理]
设 $F: \mathcal{C}^{op} \to \mathbf{Set}$ 是反变函子。则：
$$\text{Nat}(\text{Hom}(-, A), F) \cong F(A)$$
这个同构是自然的。
\end{theorem}

% ================================
% 更多填充内容以达到文件大小目标
% ================================

\subsection{补充数学公式}

以下是更多常用的数学公式，用于测试大文件的渲染性能：

\begin{align}
\int_0^\infty x^n e^{-x} dx &= n! \\
\int_0^\infty e^{-x^2} dx &= \frac{\sqrt{\pi}}{2} \\
\int_0^{\pi/2} \sin^n x \, dx &= \frac{(n-1)!!}{n!!} \cdot \begin{cases} \frac{\pi}{2} & n \text{ 偶} \\ 1 & n \text{ 奇} \end{cases} \\
\int_0^\infty \frac{\sin x}{x} dx &= \frac{\pi}{2} \\
\int_0^\infty \frac{x}{e^x - 1} dx &= \frac{\pi^2}{6} \\
\int_0^1 \frac{\ln x}{1 - x} dx &= -\frac{\pi^2}{6} \\
\sum_{n=1}^\infty \frac{1}{n^s} &= \zeta(s) \quad (\text{Re}(s) > 1) \\
\prod_{p \text{ prime}} \frac{1}{1 - p^{-s}} &= \zeta(s) \\
\Gamma(z) \Gamma(1-z) &= \frac{\pi}{\sin(\pi z)} \\
\zeta(2n) &= \frac{(-1)^{n+1} B_{2n} (2\pi)^{2n}}{2(2n)!}
\end{align}

\begin{align}
\sum_{k=0}^n \binom{n}{k}^2 &= \binom{2n}{n} \\
\sum_{k=0}^n \binom{n}{k} \binom{m}{r-k} &= \binom{n+m}{r} \\
\sum_{k=0}^r \binom{n+k}{k} &= \binom{n+r+1}{r} \\
\sum_{k=0}^n (-1)^k \binom{n}{k} &= 0 \\
\sum_{k=0}^n k \binom{n}{k} &= n \cdot 2^{n-1} \\
\sum_{k=0}^n k^2 \binom{n}{k} &= n(n+1) \cdot 2^{n-2}
\end{align}

更多矩阵运算公式：

\begin{align}
(AB)^{-1} &= B^{-1} A^{-1} \\
(A^T)^{-1} &= (A^{-1})^T \\
\det(AB) &= \det(A) \det(B) \\
\det(A^{-1}) &= (\det A)^{-1} \\
\text{tr}(AB) &= \text{tr}(BA) \\
\text{tr}(A^T) &= \text{tr}(A) \\
e^{A+B} &= e^A e^B \quad \text{if } [A, B] = 0 \\
\frac{d}{dt} e^{At} &= A e^{At} = e^{At} A
\end{align}

向量微积分公式：

\begin{align}
\nabla(fg) &= f \nabla g + g \nabla f \\
\nabla \cdot (f\mathbf{A}) &= f \nabla \cdot \mathbf{A} + \mathbf{A} \cdot \nabla f \\
\nabla \times (f\mathbf{A}) &= f \nabla \times \mathbf{A} + \nabla f \times \mathbf{A} \\
\nabla \cdot (\mathbf{A} \times \mathbf{B}) &= \mathbf{B} \cdot (\nabla \times \mathbf{A}) - \mathbf{A} \cdot (\nabla \times \mathbf{B}) \\
\nabla \times (\nabla \times \mathbf{A}) &= \nabla(\nabla \cdot \mathbf{A}) - \nabla^2 \mathbf{A} \\
\nabla \times (\nabla f) &= \mathbf{0} \\
\nabla \cdot (\nabla \times \mathbf{A}) &= 0
\end{align}

积分定理：

\begin{theorem}[Stokes 定理]
$$\oint_{\partial S} \mathbf{F} \cdot d\mathbf{r} = \iint_S (\nabla \times \mathbf{F}) \cdot d\mathbf{S}$$
\end{theorem}

\begin{theorem}[Gauss 散度定理]
$$\iiint_V \nabla \cdot \mathbf{F} \, dV = \oiint_{\partial V} \mathbf{F} \cdot d\mathbf{S}$$
\end{theorem}

\begin{theorem}[Green 定理]
$$\oint_C (P \, dx + Q \, dy) = \iint_D \left( \frac{\partial Q}{\partial x} - \frac{\partial P}{\partial y} \right) dA$$
\end{theorem}

这个超大型文档到此结束。文档包含了大量的数学定理、公式和证明，涵盖了数学的多个核心领域。通过这个文档的渲染测试，可以全面评估 LaTeX 编辑器在处理复杂大型文档时的性能表现。


% ================================
% 最终补充内容
% ================================

\section{附加材料}

\subsection{数论基础}

\begin{theorem}[算术基本定理]
每个大于 1 的正整数都可以唯一地表示为素数的乘积（不计顺序）。
$$n = p_1^{a_1} p_2^{a_2} \cdots p_k^{a_k}$$
\end{theorem}

\begin{theorem}[欧拉定理]
若 $\gcd(a, n) = 1$，则 $a^{\phi(n)} \equiv 1 \pmod{n}$，其中 $\phi(n)$ 是欧拉函数。
\end{theorem}

\begin{theorem}[费马小定理]
若 $p$ 是素数，$\gcd(a, p) = 1$，则 $a^{p-1} \equiv 1 \pmod{p}$。
\end{theorem}

\begin{theorem}[中国剩余定理]
设 $m_1, \ldots, m_k$ 两两互素，则同余方程组
$$x \equiv a_i \pmod{m_i}, \quad i = 1, \ldots, k$$
在模 $m_1 \cdots m_k$ 下有唯一解。
\end{theorem}

\begin{theorem}[二次互反律]
设 $p, q$ 是不同的奇素数，则
$$\left(\frac{p}{q}\right) \left(\frac{q}{p}\right) = (-1)^{\frac{p-1}{2} \cdot \frac{q-1}{2}}$$
其中 $\left(\frac{\cdot}{\cdot}\right)$ 是 Legendre 符号。
\end{theorem}

\begin{theorem}[素数定理]
$$\pi(x) \sim \frac{x}{\ln x}$$
其中 $\pi(x)$ 是不超过 $x$ 的素数个数。
\end{theorem}

\subsection{抽象代数}

\begin{theorem}[同构定理]
\textbf{第一同构定理：} 若 $\phi: G \to H$ 是群同态，则
$$G / \ker \phi \cong \text{Im } \phi$$

\textbf{第二同构定理：} 若 $H \leq G$，$N \trianglelefteq G$，则
$$HN / N \cong H / (H \cap N)$$

\textbf{第三同构定理：} 若 $K \trianglelefteq H \trianglelefteq G$，$K \trianglelefteq G$，则
$$(G/K) / (H/K) \cong G/H$$
\end{theorem}

\begin{theorem}[Sylow 定理]
设 $G$ 是有限群，$|G| = p^a m$，$\gcd(p, m) = 1$。
\begin{enumerate}
\item $G$ 有 $p^a$ 阶子群（Sylow $p$-子群）
\item 所有 Sylow $p$-子群共轭
\item Sylow $p$-子群的个数 $n_p \equiv 1 \pmod{p}$ 且 $n_p | m$
\end{enumerate}
\end{theorem}

\begin{theorem}[有限生成 Abel 群的结构定理]
每个有限生成 Abel 群同构于
$$\mathbb{Z}^r \oplus \mathbb{Z}_{n_1} \oplus \mathbb{Z}_{n_2} \oplus \cdots \oplus \mathbb{Z}_{n_k}$$
其中 $n_1 | n_2 | \cdots | n_k$。
\end{theorem}

\subsection{场论}

\begin{definition}[域扩张]
若 $F \subseteq K$ 是域，称 $K$ 是 $F$ 的扩张，记为 $K/F$。

扩张次数：$[K:F] = \dim_F K$
\end{definition}

\begin{theorem}[塔公式]
若 $F \subseteq E \subseteq K$，则
$$[K:F] = [K:E][E:F]$$
\end{theorem}

\begin{theorem}[Galois 基本定理]
设 $K/F$ 是有限 Galois 扩张，$G = \text{Gal}(K/F)$。则存在双射
$$\{\text{中间域 } E\} \leftrightarrow \{\text{子群 } H \leq G\}$$
给出 $E \mapsto \text{Gal}(K/E)$，$H \mapsto K^H$（$H$ 的固定域）。

且 $[K:E] = |H|$，$[E:F] = [G:H]$。
\end{theorem}

\begin{theorem}[代数闭包]
每个域 $F$ 都有代数闭包 $\bar{F}$，且在同构意义下唯一。
\end{theorem}

\subsection{实分析补充}

\begin{theorem}[Weierstrass 逼近定理]
在 $[a, b]$ 上的每个连续函数都可以被多项式一致逼近。
\end{theorem}

\begin{theorem}[Stone-Weierstrass 定理]
设 $X$ 是紧豪斯多夫空间，$\mathcal{A} \subseteq C(X, \mathbb{R})$ 是子代数，分离点且包含常函数。则 $\mathcal{A}$ 在 $C(X, \mathbb{R})$ 中稠密（关于一致收敛）。
\end{theorem}

\begin{theorem}[Arzelà-Ascoli 定理]
设 $K$ 是紧度量空间，$\mathcal{F} \subseteq C(K)$。$\mathcal{F}$ 相对紧当且仅当：
\begin{enumerate}
\item $\mathcal{F}$ 一致有界
\item $\mathcal{F}$ 等度连续
\end{enumerate}
\end{theorem}

\begin{theorem}[Radon-Nikodym 定理]
设 $\mu, \nu$ 是 $\sigma$-有限测度，$\nu \ll \mu$（$\nu$ 关于 $\mu$ 绝对连续）。则存在可测函数 $f \geq 0$ 使得
$$\nu(E) = \int_E f \, d\mu$$
$f$ 称为 Radon-Nikodym 导数，记为 $\frac{d\nu}{d\mu}$。
\end{theorem}

\begin{theorem}[Fubini 定理]
设 $(X, \mu), (Y, \nu)$ 是 $\sigma$-有限测度空间，$f \in L^1(X \times Y)$。则
$$\int_{X \times Y} f \, d(\mu \times \nu) = \int_X \left( \int_Y f(x, y) \, d\nu(y) \right) d\mu(x) = \int_Y \left( \int_X f(x, y) \, d\mu(x) \right) d\nu(y)$$
\end{theorem}

\subsection{更多复分析}

\begin{theorem}[Riemann 映射定理]
$\mathbb{C}$ 的每个单连通真子区域都共形等价于单位圆盘。
\end{theorem}

\begin{theorem}[Picard 大定理]
若 $f$ 在 $z_0$ 处有本质奇点，则在 $z_0$ 的任意去心邻域内，$f$ 取遍 $\mathbb{C}$ 中所有值，至多除一个例外。
\end{theorem}

\begin{theorem}[辐角原理]
设 $f$ 在闭曲线 $\gamma$ 内除有限个极点外解析。则
$$\frac{1}{2\pi i} \oint_\gamma \frac{f'(z)}{f(z)} dz = Z - P$$
其中 $Z$ 是零点个数（计重数），$P$ 是极点个数（计阶数）。
\end{theorem}

\begin{theorem}[Rouché 定理]
设 $f, g$ 在闭曲线 $\gamma$ 内解析，在 $\gamma$ 上连续，且在 $\gamma$ 上 $|f| > |g|$。则 $f$ 和 $f + g$ 在 $\gamma$ 内有相同个数的零点。
\end{theorem}

\subsection{附加练习题}

以下是一些综合练习题：

\begin{enumerate}
\item 证明：若 $f$ 在 $\mathbb{R}$ 上连续，且对所有 $x$ 有 $f(x + y) = f(x) + f(y)$，则 $f(x) = cx$ 对某常数 $c$。

\item 设 $A$ 是 $n \times n$ 实矩阵，满足 $A^2 = A$。证明 $\text{tr}(A) = \text{rank}(A)$。

\item 证明：$\sum_{k=1}^{n} \frac{1}{k}$ 不是整数（$n > 1$）。

\item 设 $f: [0, 1] \to [0, 1]$ 连续。证明存在 $x \in [0, 1]$ 使得 $f(x) = x$（Brouwer 不动点定理的一维情形）。

\item 计算积分 $\int_0^\infty \frac{x^3}{e^x - 1} dx$。

\item 证明：若 $p$ 是奇素数，则 $x^2 \equiv -1 \pmod{p}$ 有解当且仅当 $p \equiv 1 \pmod{4}$。

\item 设 $G$ 是阶为 $p^2$（$p$ 素数）的群。证明 $G$ 是 Abel 群。

\item 证明：$\mathbb{Q}(\sqrt{2}, \sqrt{3}) = \mathbb{Q}(\sqrt{2} + \sqrt{3})$。

\item 设 $f$ 是整函数，$|f(z)| \leq M|z|^n$ 对所有 $z$ 成立。证明 $f$ 是次数不超过 $n$ 的多项式。

\item 证明 Cantor 定理：对于任意集合 $A$，$|A| < |\mathcal{P}(A)|$。
\end{enumerate}

这个文档现在包含了足够多的数学内容，可以用于测试 LaTeX 编辑器在处理大型复杂文档时的性能。

\end{document}
